\documentclass[11pt,a4paper]{article}
\usepackage[margin=26mm]{geometry}
\usepackage{amsmath,amssymb,amsthm,booktabs}
\usepackage[hidelinks]{hyperref}
\newtheorem{theorem}{Theorem}
\newtheorem{proposition}[theorem]{Proposition}
\newtheorem{remark}[theorem]{Remark}
\newcommand{\R}{\mathbb{R}}
\newcommand{\sgn}{\operatorname{sgn}}
\newcommand{\ind}{\mathbf{1}}
\title{Gaussian Geometry and Continuous Median Depths:\\
Regularity, Adjoint Consistency, and Numerical Stability}
\author{Working technical note}
\date{16 October 2026}
\begin{document}
\maketitle
\begin{abstract}
We audit a finite one-dimensional Gaussian stochastic-solid ray model.
We characterize finite quantile roots and their singular parameter configurations,
derive a square-root normal form for coincident-peak mixtures, distinguish
differentiability of a median from that of its interface slope, and construct
weighted-quantile root brackets with an explicit certificate.
An exact identity separates ray transmittance from the associated vacancy product.
We also distinguish depth and Jacobian convergence in a thin-layer limit,
including regular finite-thickness counterexamples to naive gradient convergence.
An interval-quantile mean has a slope-free, continuously differentiable spatial
integral and bounded conditional gradients, at the cost of mixing surfaces
when its quantile band spans a transmittance plateau.
These are mathematical supplements and implementation diagnostics.
For an ideal world-space Gaussian model, we characterize when summed
absolute directional attenuation admits a single-normal representation,
and give a sharp spectral lower bound with raw-parameter interval witnesses.
Matched-primitive contrasts motivate a conditional reparameterization
recovering exact pinhole ray coefficients from existing affine quantities.
We also give conditional local Jacobian error bounds for adjoint refinement.
Separate preprocessing audits diagnose chain-rule discrepancies, a
quaternion orthogonality defect with sharp quadratic axis-ratio order,
a tight conditional scale-VJP defect formula with a quartic relative-error
counterexample, and cancellation in projected covariance determinants.
Novelty, reconstruction improvements, and end-to-end acceleration have not been established.
\end{abstract}

\section{Scope and model}
Geometry-Grounded Gaussian Splatting (GGGS) uses continuous transmittance and
implicit median-depth differentiation~\cite{gggs}. This note analyzes its
one-dimensional branch model with a fixed, finite contributor list.
Let $N\geq1$, $0<\alpha_i<1$, $\sigma_i>0$, $p_i\in\R$, and define
\begin{align}
 g_i(t)&=\alpha_i\exp\!\left(-\frac{(t-p_i)^2}{2\sigma_i^2}\right),&
 v_i(t)&=\sqrt{1-g_i(t)},\\
 T_i(t)&=\begin{cases}
 v_i(t),&t\leq p_i,\\
 (1-\alpha_i)/v_i(t),&t>p_i,
 \end{cases}&
 T(t)&=\prod_{i=1}^N T_i(t),\qquad V(t)=\prod_{i=1}^N v_i(t).
\end{align}
The ray origin is $-\infty$. Here $\alpha_i$ is an effective ray amplitude,
not necessarily a learnable three-dimensional center opacity.
Write
\[
 H(t)=-\log T(t),\quad a_i=-\log(1-\alpha_i),\quad
 A=\sum_i a_i,\quad h=-\log q,\qquad 0<q<1.
\]
Screen-space affine projection, contributor changes, opacity clipping, early
termination, finite precision, and the degenerate plane branch are outside the
continuous theorems below. A per-view affine vacancy product need not be the
restriction of one common scalar field in space.

\section{Root existence and regularity}
\begin{theorem}[Finite roots and the complete critical configuration]
\label{thm:root}
$H$ is continuously differentiable and strictly increasing on $\R$, with limits
$0$ and $A$. The equation $T(t)=q$ has a unique finite root if and only if $A>h$.
At a finite root $t^*$, $H_t(t^*)=0$ if and only if all peaks equal $t^*$.
Equivalently, a critical root exists precisely when
\[
 p_1=\cdots=p_N=p,\qquad A=2h
 \quad\left(\prod_i(1-\alpha_i)=q^2\right).
\]
At every noncritical root, the root is locally $C^1$ in the continuous
parameters while the contributor list is fixed.
\end{theorem}
\begin{proof}
Direct differentiation of either branch, including their common limiting
derivative at $p_i$, gives
\[
 H_t(t)=\sum_i
 \frac{g_i(t)|t-p_i|}{2\sigma_i^2(1-g_i(t))}.
\]
Each term is nonnegative and vanishes only at its own peak. Every nonempty
interval therefore has a strictly positive integral of $H_t$, proving strict
monotonicity. The limits follow from the branch definitions, so the intermediate
value theorem proves the root criterion. The sum vanishes precisely when every
peak equals $t$; at such a point $H(p)=A/2$. The branch values and first
parameter derivatives match at a peak. The implicit function theorem then
applies wherever $H_t>0$.
\end{proof}

\begin{proposition}[Square-root normal form]
At a critical configuration of Theorem~\ref{thm:root}, let
\[
 K=\sum_i\frac{\alpha_i}{4\sigma_i^2(1-\alpha_i)}>0.
\]
Perturb only $\alpha_j$ to $\alpha_j+\epsilon$. For sufficiently small
nonzero $\epsilon$, the finite root satisfies
\[
 t^*(\epsilon)-p=
 -\sgn(\epsilon)\sqrt{\frac{|\epsilon|}{2(1-\alpha_j)K}}
 +O(|\epsilon|^{3/2}).
\]
In particular, the parameter-to-root map is not locally Lipschitz there.
\end{proposition}
\begin{proof}
Taylor expansion on the two sides of the common peak yields
\[
 H(p+x)=A/2+Kx|x|+O(x^4).
\]
Under the perturbation,
$A(\epsilon)/2-h=\epsilon/[2(1-\alpha_j)]+O(\epsilon^2)$ and
$K(\epsilon)=K+O(\epsilon)$. Strict monotonicity determines the sign of the
root displacement. Substitution in $H(p+x)=h$ first gives
$x^2=|\epsilon|/[2(1-\alpha_j)K]+O(\epsilon^2)$ and then the stated expansion.
The square-root leading term excludes a local Lipschitz bound.
\end{proof}

For one primitive and $q=1/2$, the critical amplitude is $\alpha=3/4$ and the
leading coefficient is $\sigma\sqrt{8/3}$. This is a singular but admissible
finite root, not evidence that such configurations explain a reconstruction
failure. Distinct peaks ensure regularity but do not provide a uniform
conditioning bound as their separation tends to zero.

\section{Thin layers: depth convergence does not imply gradient convergence}
Fix distinct ordered peaks $p_1<\cdots<p_N$, fixed $0<\alpha_i<1$,
and widths $\sigma_i=\varepsilon s_i$ with $s_i>0$ and $\varepsilon\downarrow0$.
Write $\beta_i=1-\alpha_i$, $C_0=1$, $C_i=\prod_{k\le i}\beta_k$.
All per-ray opacities, peak order and projections are held fixed.
Away from the peaks the transmittance tends to the discrete plateau
$C_i$ between $p_i$ and $p_{i+1}$.
This is a ray-level comparison with discrete median selection,
not an equivalence of full GGGS, 2DGS~\cite{twodgs}, or triangle renderers.

\begin{theorem}[Strict-crossing thin-layer limit]
Suppose $C_j<q<C_{j-1}$, and put $r=q/C_{j-1}$, $\beta=\beta_j$.
Define
\[
 z_0=\begin{cases}
 -\sqrt{-2\log((1-r^2)/\alpha_j)},&\sqrt\beta\le r<1,\\
 +\sqrt{-2\log((1-(\beta/r)^2)/\alpha_j)},&\beta<r\le\sqrt\beta.
 \end{cases}
\]
Then $(t_\varepsilon-p_j)/(\varepsilon s_j)\to z_0$, hence
$t_\varepsilon\to p_j$.
If also $q\ne C_{j-1}\sqrt{\beta_j}$, then $z_0\ne0$ and, for $N\ge2$,
\[
 t_\varepsilon=p_j+\varepsilon s_jz_0+
 O(\varepsilon e^{-c/\varepsilon^2}),\qquad
 \frac{\partial t_\varepsilon}{\partial p_i}
 =\begin{cases}
 1+O(\varepsilon^{-1}e^{-c/\varepsilon^2}),&i=j,\\
 O(\varepsilon^{-1}e^{-c/\varepsilon^2}),&i\ne j,
 \end{cases}
\]
where one may take $c=\Delta^2/(8s_{\max}^2)>0$ with minimum peak
separation $\Delta$ and $s_{\max}=\max_i s_i$.
For $N=1$ the isolated-root expression and position derivative one are exact.
Constants depend on the fixed input; the claim is not uniform as the
plateau or internal peak margin vanishes.
\end{theorem}
\begin{proof}
In a bounded scaled window $t=p_j+\varepsilon s_jz$, every other peak
is at distance at least $\Delta/2$ for sufficiently small $\varepsilon$.
Its Gaussian tail is $O(e^{-c/\varepsilon^2})$.
Previous layers therefore contribute $\beta_i$ and later layers one,
up to this error, uniformly in that window. Thus
$T_\varepsilon(t)=C_{j-1}\psi_{\alpha_j}(z)
[1+O(e^{-c/\varepsilon^2})]$, where $\psi$ is the isolated-layer
transmittance. Its strict monotonicity and $\beta<r<1$ bracket a unique
finite $z_0$ and prove scaled-root convergence.
When $z_0\ne0$, $\psi'(z_0)\ne0$, so its local inverse is Lipschitz
and the scaled-root error is $O(e^{-c/\varepsilon^2})$.
For $F=-\log T$, each layer has nonnegative slope
\[
 f_i'(t)=\frac{|t-p_i|\alpha_i
 e^{-(t-p_i)^2/(2\varepsilon^2s_i^2)}}
 {2\varepsilon^2s_i^2[1-\alpha_i e^{-(t-p_i)^2/(2\varepsilon^2s_i^2)}]}.
\]
The exact position derivative is $f_i'/\sum_k f_k'$ and these
weights sum to one. The selected slope is bounded below by a
positive constant times $\varepsilon^{-1}$; other slopes are
$O(\varepsilon^{-2}e^{-c/\varepsilon^2})$. Their ratios give the claim.
\end{proof}

The plateau margin $\min(C_{j-1}-q,q-C_j)>0$ therefore controls layer
selection, while the separate margin
$|q-C_{j-1}\sqrt{\beta_j}|>0$ controls this Jacobian limit.
The following counterexample has a positive finite-thickness slope,
so the discrepancy is not caused by applying IFT at an undefined derivative.

\begin{proposition}[Regular roots with a nonlocal Jacobian limit]
Let $p=(-1,0,1)$, $\alpha=(1/5,39/64,1/5)$, $s=(1,1,1)$ and $q=1/2$.
For every $\varepsilon>0$, the unique root is $t_\varepsilon=0$ and
\[
 \left(\frac{\partial t_\varepsilon}{\partial p_1},
 \frac{\partial t_\varepsilon}{\partial p_2},
 \frac{\partial t_\varepsilon}{\partial p_3}\right)=(1/2,0,1/2),
 \qquad
 \frac{\partial t_\varepsilon}{\partial\alpha_2}
 =-\frac{32}{5}\varepsilon^2
 \left[e^{1/(2\varepsilon^2)}-\frac15\right].
\]
Thus depth converges to the middle peak but its position Jacobian
never becomes the middle-layer selector, and opacity sensitivity diverges.
\end{proposition}
\begin{proof}
Put $E=e^{-1/(2\varepsilon^2)}$. At zero the outer vacancy factors
cancel, giving $T(0)=(4/5)(5/8)=1/2$.
The cumulative plateaus obey $C_2=5/16<q<C_1=4/5$, while
$C_1\sqrt{\beta_2}=q$. The middle slope is zero, the outer slopes equal,
and $F'(0)=E/[5\varepsilon^2(1-E/5)]>0$.
The position weights follow. Since $F_{\alpha_2}(0)=1/(2\beta_2)=32/25$,
the regular-root IFT gives the stated opacity derivative.
\end{proof}

If the plateau margin itself vanishes, even layer-valued depth
convergence can fail. With two layers of opacity $1/2$, $q=1/2$,
and $p_1<p_2$, the root for every $\varepsilon>0$ is
\[
 t_\varepsilon=\frac{s_1p_2+s_2p_1}{s_1+s_2}.
\]
Indeed, between the peaks $T=q$ is equivalent to equal squared
normalized distances to the two peaks. The root remains strictly
between them. A discrete rule's equality convention cannot be
inferred by simply taking zero thickness.

Exact symbolic checks, seven 200-digit critical cases, six plateau
cases and twenty strict noncritical cases are recorded in
\texttt{thin\_quantile\_limit\_v2.json}. Independent product-$T$
re-solving verifies opacity finite differences to $2.954\times10^{-8}$
relative error; steps shrink with the squared Gaussian tail to remain
inside the shrinking linear neighborhood. Nine input exclusions based
only on predeclared analytic margins are retained. These tests concern
an ideal finite-ray model; masks, finite subdivisions, float32 rounding,
gradient floors, full three-dimensional thinning and quality are outside
the theorem. Novelty has not been established.

\section{A quantile interval: regularity and the cost of averaging}
Interval-quantile means are established statistical functionals; see, for
example, range value at risk~\cite{rvar}. We claim no new statistic or
statistical robustness theorem. The following result concerns parameter
regularity in the present finite-ray model.

\begin{theorem}[A continuously differentiable interval depth]
Fix $C_N<q_l<q_u<1$, $h=q_u-q_l$, and define
$M=h^{-1}\int_{q_l}^{q_u}t(q)\,dq$.
In a parameter neighborhood with common fixed finite bracket $[a,b]$
such that $T(a)>q_u$ and $T(b)<q_l$, one has
\[
 M=a+\int_a^b\operatorname{clip}\left(\frac{T(t)-q_l}{h},0,1\right)dt,
 \qquad
 M_\theta=\frac1h\int_{\{q_l<T(t)<q_u\}}T_\theta(t)\,dt.
\]
The map $M$ is $C^1$ in the internal continuous parameters, even if an
individual quantile inside the band is critical. For the fixed peak and
opacity parameterization,
\[
 0\le M_{p_i}\le1,\qquad \sum_i M_{p_i}=1,\qquad
 -\frac{b-a}{h(1-\alpha_i)}\le M_{\alpha_i}\le0.
\]
These bounds require no positive lower bound on the root slope.
\end{theorem}
\begin{proof}
For fixed $q$, $t(q)-a$ is the length of
$\{t\in[a,b]:q<T(t)\}$. Fubini's theorem gives the clipped integral.
On a common compact parameter neighborhood, $T$ and its parameter
partials are continuous and bounded on $[a,b]$. The clip map is
Lipschitz, and its exceptional levels $T=q_l,q_u$ each contain only one
point by strict monotonicity. Dominated convergence therefore permits
parameter differentiation and proves continuity of the integral gradient.
This argument does not differentiate a critical individual root.
Since $T_{p_i}=Tf_i'\ge0$ and $\sum_iT_{p_i}=-T_t$, integration from
$t(q_u)$ to $t(q_l)$ gives the position identities.
For $e=\exp(-(t-p_i)^2/(2\sigma_i^2))$, the two opacity log-partials are
\[
 -\frac{e}{2(1-\alpha_i e)},\qquad
 -\frac1{1-\alpha_i}+\frac{e}{2(1-\alpha_i e)}.
\]
Both are nonpositive with magnitude at most $(1-\alpha_i)^{-1}$.
Use $T\le1$ and integrate to obtain the opacity bound.
\end{proof}

The bound is local in a common bracket and depends on band width and
distance from $\alpha_i=1$. It is not uniform as $h\downarrow0$ or
over unbounded depth ranges, and is not a bound on all world parameters.

\begin{proposition}[Thin interval depths can mix surfaces]
Under the fixed, distinct-peak thin-layer assumptions above, put
$\ell_i=|[q_l,q_u]\cap[C_i,C_{i-1}]|$. Then
\[
 M_\varepsilon\longrightarrow\frac1h\sum_i p_i\ell_i,
 \qquad (M_\varepsilon)_{p_i}\longrightarrow\frac{\ell_i}{h}.
\]
If no internal plateau $C_i$, $i<N$, equals either band endpoint, then
\[
 (M_\varepsilon)_{\alpha_l}\longrightarrow
 -\frac1{h(1-\alpha_l)}
 \sum_{\substack{l\le i<N\\q_l<C_i<q_u}}
 (p_{i+1}-p_i)C_i.
\]
\end{proposition}
\begin{proof}
The quantile roots have a common finite bracket for
$0<\varepsilon\le\varepsilon_0$: Gaussian tails outside a sufficiently
large interval are uniformly small. Almost every $q$ in the fixed band
has a strict crossing and a noncritical isolated-layer root, so the
previous depth and position-weight limits apply. Depths are uniformly
bounded and the position weights lie in $[0,1]$; dominated convergence
over $q$ gives the first two limits. At the finitely many exceptional
levels no individual IFT derivative is needed. The spatial and quantile
gradient integrals agree by the monotone change of variables away from
those levels. For the opacity limit use the spatial formula instead:
between $p_i$ and $p_{i+1}$, $T_\varepsilon\to C_i$ and
$(T_\varepsilon)_{\alpha_l}\to-C_i/(1-\alpha_l)$ for $i\ge l$, otherwise
zero. Endpoint exclusion ensures convergence of the band indicator;
the common bracket and opacity bound permit dominated convergence.
\end{proof}

An interval entirely inside one layer crossing preserves that layer in
the thin limit, whereas an interval spanning a plateau mixes depths.
For $p=(-1,1)$ and $\alpha=(1/2,1/2)$, bands $[.45,.55]$ and
$[.40,.55]$ tend to depths $0$ and $+1/3$, respectively; neither is a
surface point. The protocol initially wrote $-1/3$ for the latter by
a sign error. The unchanged protocol, observed $+1/3$ and dated erratum
are preserved; no result-based parameter selection was made.

Seven prescribed critical-family cases, a singular single-layer median,
and twelve plateau/band cases are recorded in
\texttt{quantile\_band\_v2.json}. For the three-layer family at
$\varepsilon=.075$, the band $[.45,.55]$ has middle-opacity derivative
$-0.532260208$ versus approximately $-1.44629\times10^{37}$ for the
single median; the functions are different.
Spatial and independently re-solved quantile integrals agree in the
prescribed forward cases. The first double-precision small-component
finite difference failed; its source and log are archived.
Re-integrating at 60 digits with the same $10^{-6}$ step passes the
original $10^{-5}$ scaled gate (floor $10^{-8}$), with maximum error
$2.651\times10^{-11}$. This prototype establishes no CUDA speed,
reconstruction gain, or method novelty, and averaging must be tested
for surface bias before any training claim.

A subsequent predeclared frozen-ray development check reuses all 3072
uniformly sampled rays from three saved DTU models. Of these, 2621 have
a finite median and 2614 have the entire $[.45,.55]$ band finite.
All 2614 spatial-integral checks and 24 prescribed independent
forward/finite-difference checks pass. However, 2051/2614 (78.4621\%)
exceed the predeclared hypothetical thin-mixing threshold, so the
universal replacement development gate fails and no training is started.
This proxy is a fixed-parameter thinning calculation, not a measured
surface error. The band-to-median opacity-gradient norm ratio has
median $1.01412$ and maximum $267.366$; averaging does not reduce every
ray's sensitivity. Width-normalized depth shifts exceed $.1$ on
204/2614 rays, with maximum $145.987$.
All coefficients are saved as ragged NPZ arrays with model/source hashes.
These are correlated, frozen affine-reference rays, not independent
training repeats, full CUDA derivatives, or reconstruction evidence.
Two subsequently selected extreme rays pass independent product-$T$
full forward finite differences (maximum relative position/opacity
errors $6.50\times10^{-8}$/$1.60\times10^{-10}$).
A separate $q$-space full-VJP quadrature did not converge; its failed
source and log are preserved. The forward-difference check does not
repair that quadrature or reverse the failed development gate.

\section{Summed directional attenuation and a single-normal field}

This section concerns ideal world-space Gaussians, separately from the
per-view affine ray model used elsewhere in this note. Let
\[
 G_i(x)=a_i\exp\!\left(-\tfrac12(x-\mu_i)^TP_i(x-\mu_i)\right),
 \qquad 0<a_i<1,\quad P_i\succ0,
\]
and define
\begin{equation}\label{eq:directional-world}
 v_i=\sqrt{1-G_i},\qquad
 b_i=\nabla\log v_i=\frac{G_i}{2(1-G_i)}P_i(x-\mu_i),\qquad
 K_x(w)=\sum_i|w^Tb_i|.
\end{equation}
Fix $x$ throughout the following statements. The question is whether
$K_x(w)$ equals $|w^Tb|$ for some single vector $b$ and every unit $w$.
This is the single-normal absolute-gradient form. It is more restrictive
than the general normal distributions allowed by Objects as
Volumes~\cite{oav}. GGGS already states a near-side equivalence
condition~\cite{gggs}, and Gaussian Wrapping already discusses
multi-view limitations of non-oriented attenuation~\cite{wrapping}.
We do not claim the general issue as a new observation.

\begin{theorem}[Pointwise compatibility]\label{thm:directional-rank}
There exists $b\in\R^3$ such that $K_x(w)=|w^Tb|$ for every unit $w$
if and only if all nonzero $b_i$ are collinear. If the candidate is
prescribed to be the product-vacancy gradient
$b=\nabla\log\prod_i v_i=\sum_i b_i$, equality holds if and only if
all nonzero $b_i$ lie on the same half-line.
\end{theorem}
\begin{proof}
If $b=0$, the sum of nonnegative terms vanishes for every $w$, so
every $b_i=0$. Otherwise, for every unit $w\perp b$, equality implies
$\sum_i|w^Tb_i|=0$. Each $b_i$ is therefore perpendicular to the
entire plane $b^\perp$, hence parallel to $b$. Conversely write
$b_i=c_i u$ with $\|u\|=1$ and choose $b=(\sum_i|c_i|)u$.
For the prescribed product gradient the remaining condition is
$\sum_i|c_i|=|\sum_i c_i|$, precisely the common-sign condition.
The all-zero case satisfies both statements vacuously.
\end{proof}

The evenness $K_x(-w)=K_x(w)$ always holds, but does not imply this
single-normal representation. Pointwise existence of $b$ also does
not by itself establish a globally integrable gradient field.
For context, $K_x$ is the support function of the centered zonotope
$\sum_i[-b_i,b_i]$, whereas $|w^Tb|$ is that of the segment $[-b,b]$.
These support-function representations and the broader zonotope
approximation problem are established convex geometry~\cite{siegel,torres}.

\begin{theorem}[Sharp spectral obstruction]\label{thm:directional-bound}
Let $S=\sum_i b_i b_i^T$ have eigenvalues
$\lambda_1\ge\lambda_2\ge\lambda_3\ge0$. Then
\begin{equation}\label{eq:directional-error}
 E_x:=\inf_{b\in\R^3}\sup_{\|w\|=1}
       \bigl|K_x(w)-|w^Tb|\bigr|\ \ge\ \sqrt{\lambda_2(S)}.
\end{equation}
The constant in this universal inequality is attainable by valid Gaussians.
\end{theorem}
\begin{proof}
For any nonzero candidate $b$, the two-dimensional plane $b^\perp$
intersects the span of two orthonormal eigenvectors for the largest
two eigenvalues in a nonzero vector. Choose a unit $w$ in that
intersection. Then $w^Tb=0$ and
\[
 \bigl|K_x(w)-|w^Tb|\bigr|=K_x(w)
 \ge\sqrt{\sum_i(w^Tb_i)^2}
 =\sqrt{w^TSw}\ge\sqrt{\lambda_2(S)}.
\]
For $b=0$ choose an eigenvector for $\lambda_1$ instead. Taking the
infimum proves the bound. For sharpness, take $b_1=e_1,b_2=e_2$;
then $\lambda_2=1$, and the candidate $b=e_1$ has uniform error
$\sup_{\|w\|=1}|w_2|=1$. To realize these vectors by Gaussians, put
\[
 a_1=a_2=\tfrac12,\quad
 s=\frac{e^{-1/2}}{4-2e^{-1/2}}>0,\quad
 P_1=P_2=s^{-2}I,\quad \mu_1=-s e_1,\quad \mu_2=-s e_2.
\]
Equation~\eqref{eq:directional-world} gives $b_1=e_1,b_2=e_2$ at
$x=0$. Replacing $\mu_2$ by $+s e_1$ instead gives $b_1=-b_2=e_1$:
a single-vector representation exists, but the product gradient is zero.
\end{proof}

The argument uses the standard matrix variational principle and
$\ell_1\ge\ell_2$, not a new matrix inequality. It supplies a simple
quantitative certificate for this particular attenuation model.

\begin{proposition}[A two-term certificate]\label{prop:directional-pair}
For any two terms let $d=\|b_1\times b_2\|^2$ and
$r=\|b_1\|^2+\|b_2\|^2>0$. Then $E_x\ge\sqrt{d/r}$,
even after adding any other Gaussian terms.
\end{proposition}
\begin{proof}
The two potentially nonzero eigenvalues of
$S_{12}=b_1b_1^T+b_2b_2^T$ have product $d$ and sum $r$.
Thus $\lambda_2(S_{12})=d/\lambda_1(S_{12})\ge d/r$.
Since $S-S_{12}\succeq0$, eigenvalue monotonicity and
Theorem~\ref{thm:directional-bound} give the claim.
\end{proof}

For strictly positive interval bounds $d_-$ and $r_+$, a saved
rational $L>0$ satisfying $L^2r_+\le d_-$ is a conservative
absolute error bound, conditional on the validity of the input
enclosures. The double-precision full-model spectrum and any
normalization by $\sum_i\|b_i\|$ are separate nominal measurements.
A subset denominator cannot certify a full-model relative bound.

\paragraph{Real-parameter audit.}
We fixed 24 query points: for each of eight saved cameras in each of
scans 37, 55 and 63, the first saved finite affine-median ray was
backprojected using the saved calibration. These are neither GT
surface samples nor certified world-model median roots. Each original
PLY hash matched; all 1,010,226, 1,052,385 and 510,660 primitives,
respectively, entered the nominal world-field evaluation. Underflow
counts are retained. For each point two original raw-parameter
primitives give a strictly positive 80-digit interval certificate;
the smallest is only $3.4193\times10^{-10}$ in inverse world-coordinate
units, so positivity alone does not establish a visible effect.
Across these fixed points the nominal ratio
$\sqrt{\lambda_2(S)}/\sum_i\|b_i\|$ has minimum $0.00900015$,
median $0.06790113$ and maximum $0.21308058$. These are not
interval-certified relative bounds, surface prevalence estimates,
pixel errors or predicted quality improvements.

An independent replay verified 13 source files and their actual archive
bytes, 26 artifacts, the three PLY hashes, every query backprojection
and all 48 raw primitive records. Exact rational arithmetic verifies
all 24 saved lower-bound inequalities. Independent 120-digit scalar
differentiation of $\log v_i$ places all 144 spatial derivative
components inside both 80- and 100-digit interval enclosures.
The latter are nested inside the former. Symbolic identities and
64 predeclared numerical direction witnesses also pass.
Files \texttt{directional\_field\_theory\_v1.json} and
\texttt{directional\_field\_world\_report\_v1.json} record these checks.
The mpmath interval implementation is explicitly experimental~\cite{mpmathiv};
this is reproducible numerical evidence, not proof-assistant certification.
No production affine-renderer, reconstruction or training conclusion
follows from this world-model audit. A universal single-normal penalty
could also suppress legitimate edges or multiple layers. Novelty and
the relevance of this obstruction to measured reconstruction errors
remain unresolved.

\paragraph{Same-ray world/affine comparison.}
An additional exploratory check keeps exactly these 24 saved rays and
restricts every world Gaussian to each ray without the affine
contributor selection. For finite camera origin $o$, it uses
$T_{\mathrm{cam}}(t)=T_{-\infty}(t)/T_{-\infty}(0)$, not an
unconditional identification of the two start points. All 24 camera
medians are finite. The world-ray and separate pointwise evaluations
agree in $\log V$ within $1.98\times10^{-13}$ and in attenuation
within a scaled $8.47\times10^{-13}$. The largest root log-residual
is $3.04\times10^{-11}$; this is not an interval-certified root error.
The absolute world/affine median displacement divided by saved affine
depth has median $0.0126254\%$ over the 24 rays, with per-scan medians
$0.0087655\%$, $0.0115090\%$ and $0.0526133\%$.
The largest case is $2.69459\%$, scan37/camera0041, where depth
changes from $6.089402428$ to $6.253487084$ in world-coordinate units.
All cases, including this tail, are retained. Archive and arithmetic
replay verifies 40 sources, 26 artifacts, three PLYs and all rows;
it does not independently re-solve every full-model root.
This comparison jointly changes affine shapes, clipping, primitive
selection and opacity clamping. It isolates neither a perspective
effect nor a causal reconstruction defect. In particular, a large
directional-attenuation difference need not imply a large median-depth
shift. Without GT or cross-view error comparisons, neither result is
shown more accurate. See \texttt{world\_affine\_rays\_report\_v1.json}.

\section{Matched ray factors and a perspective reparameterization}

\paragraph{Fixed-input mechanism check.}
For the same 24 saved rays, an ID-carrying CPU projection reproduces
every old projection array and every selected ray parameter bitwise.
On the identical selected primitive set, we evaluate all eight
combinations of affine/world peak, width and peak amplitude. All
192 roots are finite; independent re-solving with the original
\texttt{log\_fields} function agrees within $2.69\times10^{-13}$.
We average marginal increments over all six factor-addition orders
(a standard Shapley-style allocation), then follow a fixed path
removing opacity capping, prefix stopping, eligibility, pixel-tile,
projection-validity and near-plane restrictions, and finally changing
the origin convention. All intermediate roots are finite. The
decomposed signed increments sum to the full world-camera/affine
contrast within $4.34\times10^{-19}$.

The medians of absolute increments divided by old depth are
$0.00209194\%$ for peak, $0.0000389273\%$ for width,
$0.000826298\%$ for amplitude and $0.0115894\%$ for removing
affine power/opacity eligibility. The median cancellation fraction
$1-|\sum_j\Delta_j|/\sum_j|\Delta_j|$ is $0.291519$.
In the largest tail case, the peak allocation is $+0.204732767$
world units, while removing eligibility contributes $-0.041186829$;
the final net is $+0.164084657$. Neither a net difference nor this
order-dependent allocation identifies a unique physical cause or GT
error. All 69 source/archive entries, 50 artifacts, three PLYs and
24 packed membership records are verified; the large selection-stage
roots are not independently re-solved by that replay.

\paragraph{Exact-real setting.}
Place the pinhole camera at the origin. Let a Gaussian have mean
$\mu=zq_0$, $q_0=(u_0,v_0,1)^T$, $z>0$, precision $P\succ0$
and covariance $\Sigma=P^{-1}$. For a pixel let
$q=q_0+\delta$, $\delta=(\delta_u,\delta_v,0)^T$, and
$d=q/\|q\|$. Define
\[
 r=q_0^Tq_0,\quad b=q_0^TPq_0,\quad
 t_0=z\sqrt r,\quad s_0=\sqrt{r/b},\quad
 J=\frac1z\begin{pmatrix}1&0&-u_0\\0&1&-v_0\end{pmatrix}.
\]
The center projection is unclipped and the screen covariance has no
extra kernel. World-space filtering can already be included in $\Sigma$.
Let $C=J\Sigma J^T$, $E_A=\delta_{xy}^TC^{-1}\delta_{xy}$, and
let the affine peak shift be the linear Taylor term
\begin{equation}\label{eq:affine-peak-shift}
 \Delta_A=z\left(\frac{q_0^T\delta}{\sqrt r}
       -\frac{\sqrt r\,\delta^TPq_0}{b}\right).
\end{equation}
Pixel focal-length scaling can equivalently be absorbed into $C$ and
$\delta_{xy}$. $E_A$ is computed from the conic, before opacity capping.

\begin{theorem}[Conditional affine-packet identity]\label{thm:packet}
Set
\[
 g=\frac{q_0^T\delta}{r},\quad k=\frac{\delta^T\delta}{r},\quad
 f=1+2g+k,\quad \ell=g-\frac{\Delta_A}{t_0},\quad
 \chi=E_A\left(\frac{s_0}{t_0}\right)^2,\quad
 D=(1+\ell)^2+\chi.
\]
Then $f,D>0$, and the exact one-dimensional restriction of the world
Gaussian to the unit ray $td$ has peak $p$, width $s$ and peak
amplitude $a$ given by
\begin{equation}\label{eq:packet-exact}
 p=\frac{t_0\sqrt f(1+\ell)}{D},\qquad
 s=\frac{s_0\sqrt f}{\sqrt D},\qquad
 a=a_0\exp\!\left(-\frac{E_A}{2D}\right),
\end{equation}
where $a_0$ is the Gaussian's uncapped world amplitude.
\end{theorem}
\begin{proof}
Use the basis $M=[e_1,e_2,q_0]$, with $\det M=1$ and
$JM=z^{-1}[I_2,0]$. The block inverse identity for $M^TPM$
gives the Schur complement
\[
 C^{-1}=z^2\left(P_{xy,xy}
       -\frac{(Pq_0)_{xy}(Pq_0)_{xy}^T}{b}\right),\qquad
 E_A=z^2\left(\delta^TP\delta
       -\frac{(\delta^TPq_0)^2}{b}\right).
\]
Equation~\eqref{eq:affine-peak-shift} yields
$\ell=\delta^TPq_0/b$ and
$\chi=\delta^TP\delta/b-\ell^2\ge0$.
Consequently $D=q^TPq/b>0$ and $f=\|q\|^2/r>0$.
Completing the ray quadratic gives
\[
 p=\frac{z\|q\|q^TPq_0}{q^TPq},\quad
 s^2=\frac{\|q\|^2}{q^TPq},\quad
 E_W=z^2\left(b-\frac{(q^TPq_0)^2}{q^TPq}\right)=\frac{E_A}{D}.
\]
Substitution proves all three identities.
\end{proof}

This is an algebraic expression of the same exact ray model, not a
second-order truncation. Exact perspective-correct Gaussian evaluation
already exists in HTGS~\cite{htgs} and 3DGUT~\cite{threedgut}; the
general concept is not claimed as new. Reusing an affine packet also
does not establish lower bandwidth, a speedup or a cheaper backward pass.
The exact-real identity can inherit conic roundoff or a finite-precision
covariance/precision mismatch. The linear term should be evaluated
directly as the plane-offset dot product rather than by subtracting
two already-added depths.

\paragraph{Verified domain and exclusions.}
Nine symbolic identities, including a general symmetric $3\times3$
Schur identity and the Taylor derivative, pass. All 64 predeclared
SPD cases pass the $10^{-8}$ scaled-error criterion with floor
$10^{-12}$. The real 24-ray selected sets contain 1,992 primitive--ray
inputs. Exactly 1,946 lie in the unclipped domain and all pass;
46 clipped-center inputs are retained separately and are not counted
as passes. An independent 80-digit factor-space calculation from every
applicable original PLY record gives maximum scaled errors for
$(p,s,a)$ of $(4.99\times10^{-15},1.66\times10^{-15},
2.64\times10^{-12})$. It verifies 124 actual source/archive entries,
50 artifacts, three PLY hashes and all domain classifications.
Saved binary world directions are used without silently re-normalizing
camera calibration. High precision is numerical cross-checking, not
interval or proof-assistant certification.

Files \texttt{perspective\_reparam\_theory\_v1.json} and
\texttt{perspective\_reparam\_report\_v1.json} record the checks.
The unclipped expression alone does not cover the 46 excluded inputs;
a separate clipped-reference expression follows below. Neither identity
repairs contributor selection or establishes global geometry consistency,
a full-domain renderer or a quality improvement.

\subsection{Recovering the quadratic form with a clipped projection center}

Retain $\Sigma=P^{-1}$, $P\succ0$, $z>0$ and zero screen kernel.
Let $\bar q=(\bar u,\bar v,1)^T$ be an arbitrary projection reference,
possibly clipped, while $\mu=zq_0$. Put $\bar r=\|\bar q\|^2$,
$\bar b=\bar q^TP\bar q$ and $\bar s^2=\bar r/\bar b$.
The cached radial depth is $t_0=z\|q_0\|$; write $S_2$ for the screen
conic, $h$ for the two plane coefficients, and $F$ for the focal matrix.
The existing quantities recover
\begin{equation}\label{eq:clipped-schur-recovery}
 z=\frac{t_0}{\|q_0\|},\quad \bar b=\frac{\bar r}{\bar s^2},\quad
 c=\frac{Fh\sqrt{\bar r}/z+\bar q_{xy}}{\bar r},\quad
 K=\frac{FS_2F}{z^2\bar b}.
\end{equation}
$q_0$ is obtained from the cached 2D center and calibration, without
reading the original 3D mean to help the candidate.

\begin{theorem}[Clipped-reference packet identity]\label{thm:clipped-packet}
For the zero-kernel packet evaluated at $\bar q$, $K\succ0$ and
\begin{equation}\label{eq:clipped-metric}
 \bar M^TP\bar M=\bar b
 \begin{pmatrix}K+cc^T&c\\c^T&1\end{pmatrix},\qquad
 \bar M=[e_1,e_2,\bar q].
\end{equation}
Thus $(\bar b,c,K)$ uniquely determines the six degrees of freedom of $P$.
For $q=(u,v,1)^T$, define
\[
 a=q_{xy}-\bar q_{xy},\quad m=q_{0,xy}-\bar q_{xy},\quad
 A=1+c^Ta,\quad B=1+c^Tm,\quad D=A^2+a^TKa,\quad C=AB+a^TKm,
\]
and
\begin{equation}\label{eq:clipped-gram}
 v=Am-Ba,\qquad
 \Gamma=\det(K)(a_1m_2-a_2m_1)^2+v^TKv\ge0.
\end{equation}
Then the exact unit-ray coefficients are
\[
 p=z\|q\|\frac{C}{D},\qquad s=\frac{\|q\|}{\sqrt{\bar bD}},\qquad
 \alpha=a_0\exp\!\left(-\frac{z^2\bar b\Gamma}{2D}\right).
\]
\end{theorem}
\begin{proof}
The block inverse identity in the basis $\bar M$ gives
$K=(P_{xy,xy}-\bar bcc^T)/\bar b$ with $c=(P\bar q)_{xy}/\bar b$.
The implementation's plane encoding is
$Fh=z(\bar r c-\bar q_{xy})/\sqrt{\bar r}$, yielding
\eqref{eq:clipped-schur-recovery} and \eqref{eq:clipped-metric}.
Conversely multiply by $\bar M^{-T}$ and $\bar M^{-1}$ to recover $P$.
This is the standard Schur/LDL block parameterization, not a new decomposition.

The coordinates of $q,q_0$ in this basis are $(a,1)$ and $(m,1)$.
Hence $q^TPq=\bar bD>0$, $q^TPq_0=\bar bC$ and
$q_0^TPq_0=\bar b(B^2+m^TKm)$. Expanding gives
\[
 D(B^2+m^TKm)-C^2
 =\det(K)(a_1m_2-a_2m_1)^2+(Am-Ba)^TK(Am-Ba)=\Gamma.
\]
Completing the ray quadratic proves the formulas. The nonnegative
expression avoids directly subtracting two large energy terms,
but determinants and cross terms retain floating-point sensitivity.
\end{proof}

The exact composite ray mapping is independent of $\bar q$.
An intermediate clipping operation therefore does not by itself
establish nonsmoothness of the exact final mapping. No boundary
derivative or universal FP32 guarantee follows from this observation.

\paragraph{Coverage and precision failure.}
Six general symbolic identities and 128 fixed SPD cases pass, with
112 clipped cases. All 1,992 original selected primitive--ray inputs
pass, including the 46 previously excluded, without fallback or
reselection. Independent 80-digit original-PLY evaluation covers all
1,992 inputs, with maximum scaled errors
$(4.72\times10^{-15},1.52\times10^{-15},2.61\times10^{-12})$.
Replay verifies 180 source/archive entries and 50 artifacts.
This covers the selected CPU scope, not every primitive/pixel pair
or nonzero screen kernels.

A subsequent predeclared precision screen preserves a negative result.
All 1,992 real packets pass after input rounding to FP32, which does
not include actual FP32 preprocessing. Raw FP64 packet and direct-world
AD graphs agree for all 71,712 components from 12 parameters and three
outputs; the largest scaled discrepancy is $4.19\times10^{-7}$ against
a $10^{-4}$ gate with floor $10^{-8}$. They share the raw transform
and Torch AD library; shared-primitive atomic accumulation and CUDA
backward are not tested. All sampled raw centers are away from clip boundaries.

One of 128 FP32 synthetic cases fails the $10^{-5}$ relative peak gate,
with error $5.29059\times10^{-5}$; the overall development gate stays
failed. In a post hoc 80-digit calculation, decoding the already-rounded
packet still gives peak error $1.71395\times10^{-5}$, so raising
evaluation precision alone does not repair this case. Its peak is about
$-0.02839818$ and amplitude $2.3005\times10^{-5712}$, zero by double
underflow. That limits physical relevance but is not used to reclassify
the failed gate. All inputs are retained. Precision replay verifies
233 source/archive entries, 51 artifacts and stored gradient arithmetic.
See \texttt{clipped\_packet\_report\_v1.json} and
\texttt{clipped\_packet\_precision\_report\_v1.json}.
An exploratory input-protection screen exhausts 64 combinations on
the same 1,992 real and 128 synthetic inputs. Two inclusion-minimal
combinations pass: protecting only $q_0$, or protecting the plane,
conic and camera geometry. Neither survives a locked, separately
specified test on 4,096 new same-distribution cases (seed 13017):
they fail on 13 and 11 cases, respectively, while all-double passes
all cases. The float32 baseline fails 21; rounded inputs with double
evaluation fail 11. An independent 80-digit factor oracle preserves
all classifications on the union of 24 failed cases. The two
candidates are rejected; these are synthetic numerical inputs,
not independent real-scene quality evidence.

Finite perturbation analysis explains the distinction between
absolute and relative peak sensitivity. Write $d=A^2+a^TKa$,
$C=AB+a^TKm$, $q=C/d$, and suppose
$|\Delta C|\leq e_C$, $|\Delta d|\leq e_d<d$ and
$|\Delta(z\|q_{\rm ray}\|)|\leq e_s$. Then
\[
 |\Delta q|\leq e_q:=\frac{e_C+|q|e_d}{d-e_d},\qquad
 |\Delta p|\leq |z\|q_{\rm ray}\||e_q+e_s(|q|+e_q).
\]
This follows directly from the difference of two quotients and
the product rule for finite differences. Relative sensitivity may
be large when $C$ is close to zero. Explicit bounds for the packet
variables and conditional log-transmittance/root error are recorded
in \texttt{PACKET\_CONDITIONING\_THEORY.md}; they require positive
denominator and slope margins and do not establish full-scene quality.

A two-point witness gives a limited representation obstruction.
Two legal clipped/SPD packets have identical binary32 input bytes
but their $10^{-5}$ relative peak acceptance intervals are disjoint.
A 100-digit interval check gives a separation lower bound about
$4.92775\times10^{-7}$. Any deterministic decoder receiving only
those same bytes returns the same value, so cannot satisfy both
intervals. This concerns the eight-input prototype that directly
rounds $q_0$; it excludes neither additional metadata nor all
production UV/intrinsic layouts. The interval library remains a
trusted experimental dependency, not proof-assistant certification.

A diagnostic extension also launches the copied production CUDA
preprocessing functions on all 1,992 inputs, including 46 clipped
ones, with zero culled inputs. Production-UV recovery followed by
Torch CUDA float32 decoding passes the same local gates; the largest
peak/width relative and amplitude absolute errors are
$5.53900\times10^{-7}$, $1.29782\times10^{-6}$ and
$1.18168\times10^{-5}$. All 5,976 outputs replay bitwise, and
actual CUDA denominator/determinant/energy checks report no invalid
inputs. This is a separately compiled diagnostic extension and
elementwise Torch decoder, not the installed renderer binary or a
full backward check. It does not reverse the synthetic failures.
These local checks do not establish GT quality or novelty.

The subsequent inference integration preserves legacy affine selection,
RGB/normal accumulation, prefix, initialization and subdivision rules,
and explicitly rejects backward, sampled NCC and integration calls.
Its one-image smoke passes. However, the predeclared complete-image
comparison on all 24 development views (10,834,488 pixels) rejects
the first packet implementation. All saved arrays are finite,
RGB/alpha/normal remain bitwise identical, and the packet/factor
valid-depth masks agree. Although each view's common-valid depth
error at percentile 99 is at most $1.90735\times10^{-6}$, eight
views fail the $10^{-3}$ maximum absolute-error gate: nine pixels
exceed it, with maximum $0.17964053$ in model length units.
The sparse tail is retained rather than excused by the percentile.

Both exact arms use FP64 scalar coefficient evaluation with FP32
metadata and the legacy FP32 solver. Counterbalanced complete-render
timing gives packet/factor GPU ratio geometric mean $2.6001822$
(maximum $2.6573840$); synchronized wall ratio is $2.5999135$.
This also fails the efficiency gate, and is not a training-throughput
comparison. No GT-quality or training expansion follows this version.
See \texttt{FROZEN\_EXACT\_RENDER\_RESULT.md}.

A subsequent, explicitly posthoc analysis retains all nine outliers.
It captures 17,581 tile candidates and 772 accepted terms from eight
affected views. An accidental CUDA architecture change rebuilt the
cached binaries, and restoring the old environment did not recover
their historical byte identity. The diagnostic therefore pins new
binary hashes and first requires all eight affected views' depth,
secant slope, RGB, alpha and normal outputs to replay the old saved
FP32 bit patterns. All pass. This is output equivalence, not recovery
of historical hidden buffers or instruction identity. The captured
prefix is an actual buffer; coefficients and initialization are
obtained from a separately compiled expression probe.

For both saved coefficient representations on each ray, the exact
binary-rational product $\prod_i(1-\alpha_i)$ is strictly below
$1/2$. All 18 models therefore have finite roots; the hypothesized
false-finite explanation is excluded for these samples. Independent
120-digit stable log-transmittance roots have strict bracket endpoint
signs rechecked at 180 digits. These are numerical brackets, not
outward-rounded interval certificates. At the largest outlier, in
ray-length units,
\[
 \underbrace{-0.1820801374}_{\text{observed packet minus factor}}
 =\underbrace{-0.1820989237}_{\text{represented mathematical roots}}
  +\underbrace{0.00001878635}_{\text{remaining numerical/conversion path}}.
\]
Across all nine, the absolute remainder is at most
$8.965\times10^{-5}$ ray-length units and its ratio to the absolute
observed difference is at most $2.713\%$. The remainder includes
finite subdivision, arithmetic, interpolation, probe context and
ray-to-$z$ conversion; it is not attributed solely to product rounding.
Neither FP32 representation has yet been established as the raw-world
reference. The true $-T_t$ values at the 18 roots range from
$3.1154\times10^{-4}$ to $1.9702\times10^{-2}$; these differ from
the stored secant slopes and do not supply a uniform derivative bound.

A controlled CUDA replay with the same saved coefficients and initial
bracket, eight subdivisions and five rounds matches all 18 original
depth and secant-slope bit patterns. It uses the actual initial gate
$T\leq0.45$ and search radius $0.4$. Replacing the product with FP64
stable log-transmittance leaves the large packet/factor gap present.
Thus improving the solver arithmetic alone does not eliminate the
observed representation discrepancy. The earlier replay's mistaken
$0.5$ gate is retained as an implementation erratum; all nine initial
$T$ values are below $0.393823$, so the two gates happened to agree
on these inputs. No frequency, training-quality, causal or novelty
claim follows from these selected failure cases.
See \texttt{EXACT\_OUTLIER\_RESULT.md} and
\texttt{exact\_outlier\_root\_v1.json}.

The fixed coefficient-factorial follow-up evaluates all eight choices
of packet/factor amplitude, peak and reciprocal width on each ray.
All 72 models have finite roots, with 120-digit roots and 180-digit
endpoint sign checks. Replacing only effective ray amplitude reproduces
the full replacement within $9.58250\times10^{-7}$ camera-ray units
on all nine. Replacing peak and reciprocal width together changes
the factor root by at most $9.58099\times10^{-7}$. Effective amplitude
includes off-ray energy attenuation; this does not isolate the learned
opacity parameter. The six replacement-order average is an algebraic
allocation, not a causal contribution or a jointly realizable 3D model.

\begin{proposition}[A sufficient peak-free gap bound]
Let $a<b$ and suppose no peak lies in $(a,b)$. Set
$L=\{i:p_i\leq a\}$, $R=\{i:p_i\geq b\}$,
$M_L=-\sum_{i\in L}\log(1-\alpha_i)$ and
$u_i(t)=-\tfrac12\log(1-g_i(t))$. Then, for $a\leq t\leq b$,
\[
 M_L-\sum_{i\in L}u_i(a)\ \leq\ H(t)\ \leq\
 M_L+\sum_{i\in R}u_i(b),\qquad H=-\log T.
\]
If the left bound exceeds $\log2$, a finite median root lies to
the left of $a$; if the right bound is below $\log2$, it lies to
the right of $b$.
\end{proposition}
\begin{proof}
On this interval, $H=M_L-\sum_Lu_i+\sum_Ru_i$.
Each left tail is nonnegative and decreasing, and each right tail
is nonnegative and increasing. Dropping the opposite nonnegative
tail yields the two bounds. Monotonicity of $H$ gives localization.
This uses ordinary real-arithmetic monotonicity, without a derivative
lower bound or a claim of general novelty.
\end{proof}

Taking the central 60\% of each pair's root gap, all 18 models have
no interior peak. The sufficient bounds place both models on opposite
sides in cases 1, 2 and 5; six other single-arm bounds remain
inconclusive and are retained. Independent 180-digit direct-$H$
evaluation checks 126 points. At the largest two gap midpoints,
$H_t$ is approximately $5.584\times10^{-13}$ and
$6.4\times10^{-161}$, despite regular root-local slopes.
These are high-precision numerical checks, not directed interval
certificates; a root-local derivative is not a bound across the gap.

An original-PLY reference also verifies all 772 accepted raw parameter
rows. Two unrounded high-precision objects are distinguished: a true
covariance inverse of saved FP32 activated parameters, retaining the
nonunit quaternion polynomial, and a physical covariance obtained by
high-precision activation and quaternion normalization of the raw PLY.
The saved FP32 camera affine map is inverted exactly, so $t$ denotes
unit camera-ray parameter rather than assumed unit world arclength.
Explicit covariance inversion at 80 digits and inverse-factor residual
evaluation at 140 digits agree; all 36 lower-precision root brackets
retain strict signs at higher precision. Both references find factor
closer on eight of these nine selected rays, and packet closer on
case 0. This is a reference-dependent descriptive comparison, not
GT geometry, significance or a generic precision ranking. Factor
still has about $1.20971\times10^{-3}$ ray error on case 0.

The actual common preprocessing opacity differs at the last bits
from the saved activated opacity on 96 of 772 accepted terms, with
maximum absolute difference $5.96046\times10^{-8}$. Both exact arms
share the same preprocessing opacity. A third reference holds the
activated covariance fixed and uses that common opacity; the closer
counts remain eight versus one. Thus this common amplitude difference
does not account for the large between-arm contrast. No complete
backward, quality gain, speed improvement or novelty follows from
this posthoc analysis. See
\texttt{EXACT\_OUTLIER\_COEFFICIENT\_RESULT.md}.

\subsection{A coherent factor does not remove the coefficient-precision boundary}
A fixed CPU follow-up constructs a true inverse factor and writes its
upper-triangular QR factor as six entries plus camera mean. For the
standard quaternion polynomial, with $v=(x,y,z)$,
$k=\lVert q\rVert^2-1$ and $D=1+4k\lVert v\rVert^2\ne0$,
\[
 R^{-1}=\frac{R^{\mathsf T}+4kvv^{\mathsf T}}{D}.
\]
Indeed, $Rv=v$ and
$R^{\mathsf T}R=I+4k(\lVert v\rVert^2I-vv^{\mathsf T})$ imply
$(R^{\mathsf T}+4kvv^{\mathsf T})R=DI$.
Nine symbolic polynomial components vanish exactly. The identity and
QR norm preservation are standard algebra, used here to avoid assuming
that an activated FP32 quaternion matrix is exactly orthogonal.

The preselected cache uses six FP32 triangular entries and three FP64
camera-mean entries (48 bytes per primitive). All 7,010 fixed coefficient
inputs are retained. The nine real outlier roots pass the $10^{-3}$
rule, with maximum error $4.99316\times10^{-5}$ camera-ray units.
However, 12 of the existing 4,096 geometric constructions fail the
local coefficient rules, and five of nine critical double-layer roots
fail. Thus the candidate is rejected. The 72-byte all-FP64 cache is
only a control: it still rounds the ray coefficients to FP32, and four
critical roots fail. No full-image timing or training follows.

\begin{proposition}[A finite FP32 ray-coefficient interface limitation]
There are two legal positive-definite Gaussian worlds with FP32 centers,
scales and activated amplitudes, whose correctly rounded FP32 arrays
$(p_i,\rho_i,\alpha_i)_{i=1}^2$ are identical, while their unrounded
GGGS median roots satisfy $r_A<64/125$ and $r_B>186/125$.
Consequently, every method producing the same scalar from only these
identical arrays has worst-case absolute error strictly larger than
$61/125=0.488$ in camera-ray units on the two worlds.
\end{proposition}
\begin{proof}
Use identity camera and rotation, query $(0,0,1)$, common isotropic
scale $\varepsilon=8589935/2^{32}$, peak coordinates $1/2$ and $3/2$,
and activated amplitudes $3/4$ and $1/2$. The second center has zero
transverse displacement. For the first center take respectively
\[
 y_A=7735379/2^{32},\qquad y_B=15470759/2^{33}.
\]
All values are binary32, and $y_B-y_A=2^{-33}$. The exact effective
first-layer amplitudes are
$\alpha_X=(3/4)\exp[-y_X^2/(2\varepsilon^2)]$.
Directed interval checks give
\[
 \alpha_A\simeq0.5000000205022779>1/2,\qquad
 \alpha_B\simeq0.4999999942937956<1/2.
\]
Both are strictly inside the binary32 nearest-rounding cell
$(1/2-2^{-26},1/2+2^{-25})$, so both round to $1/2$.
Peaks are exactly $1/2,3/2$, and both widths have reciprocal
$1/\varepsilon$; its correctly rounded bits are checked by exact
rational midpoint comparisons. Thus all six coefficients agree.
The represented symmetric model has exact root $1$, since the near
and far vacancy factors cancel there; finite uniqueness follows from
the positive widths and terminal transmittance $1/4<1/2$.

For the unrounded worlds, 60-digit and 120-digit directed interval
checks yield $T_A(64/125)<1/2$ and $T_B(186/125)>1/2$.
Strict monotonicity therefore gives the asserted root localization,
so $r_B-r_A>122/125$. For any common output $z$, the triangle
inequality implies
$\max(|z-r_A|,|z-r_B|)\geq|r_B-r_A|/2>61/125$.
The concrete sign/rounding certificate trusts the outward interval
implementation of mpmath; it has not been checked by a proof assistant.
\end{proof}

This existential witness was selected after observing the fixed critical
fixtures. It is not a frequency estimate or a claim about all metadata
layouts: original world parameters and higher-precision ray values
retain additional information. Complete-renderer eligibility, the
initial bracket, CUDA instructions, backward and GT geometry were not
tested by this witness. The failure is relative to the unrounded world
functional; solving the represented symmetric functional exactly still
returns its correct root $1$. The candidate result and certificate are
saved in \texttt{COHERENT\_FACTOR\_RESULT.md} and
\texttt{coherent\_factor\_audit\_v1.json}; general quantile instability
and academic novelty remain distinct questions.

\subsection{A conditional precision certificate from the original parameters}

The coefficient-alias construction motivates retaining the source geometry.
Here a cached factor can propose a root; a separate enclosure must certify it.
Let $F=\operatorname{diag}(1/s)R^{-1}A^{-1}$ be the exact saved-activated
precision factor, and let $\widehat Q,U$ be the nominal QR orthogonal and
cached triangular factors. Directed interval preprocessing supplies
$\delta_Q\geq\|\widehat Q^\top\widehat Q-I\|_2$, with $\delta_Q<1$.
The polar factor $Q$ of $\widehat Q$ is orthogonal and satisfies
\[
 \|\widehat Q-Q\|_2\leq
 \frac{\delta_Q}{1+\sqrt{1-\delta_Q}}.
\]
Indeed, the singular values of $\widehat Q$ lie between
$\sqrt{1-\delta_Q}$ and $\sqrt{1+\delta_Q}$; the largest distance to one
is bounded by the displayed expression. Thus a computable upper bound
\[
 \eta\geq\|F-\widehat Q U\|_F+
 \frac{\delta_Q\|U\|_F}{1+\sqrt{1-\delta_Q}}
\]
also bounds $\|Q^\top F-U\|_2$. No reference median is used to form it.
Let $m_*$ be the exact camera mean and let the cached mean satisfy
$\|m_*-m\|\leq\nu$. Put $M_*=Q^\top F$ and take a unit camera ray $d$.

\begin{proposition}[A residual-based coefficient enclosure]
Define $u=Ud$, $v=Um$, $\rho=\|u\|$, $p=u^\top v/\rho^2$,
$r=v-pu$, $e=\|r\|^2$, and $x=m-pd$.
Suppose $\rho>\eta$, and put
\[
 \zeta=\eta\|x\|+(\|U\|_F+\eta)\nu,
 \qquad
 D_p=\frac{\eta\sqrt e+(\rho+\eta)\zeta}{(\rho-\eta)^2},
 \qquad D_r=\zeta+D_p(\rho+\eta).
\]
The original model satisfies
\[
 |p_*-p|\leq D_p,\qquad |\rho_*-\rho|\leq\eta,
 \qquad |e_*-e|\leq2\sqrt e\,D_r+D_r^2.
\]
\end{proposition}
\begin{proof}
Write $u_*=M_*d$ and $w_*=M_*(m_*-pd)$. Then
$\|u_*-u\|\leq\eta$ and $\|w_*-r\|\leq\zeta$.
Since $u^\top r=0$,
\[
 p_*-p=\frac{u_*^\top w_*}{\|u_*\|^2},\qquad
 |u_*^\top w_*|\leq\eta\|r\|+(\rho+\eta)\zeta.
\]
The denominator is at least $(\rho-\eta)^2$.
The true perpendicular residual is $r_*=w_*-(p_*-p)u_*$, so
$\|r_*-r\|\leq D_r$. The squared-norm perturbation and reverse triangle
inequality give the remaining bounds.
\end{proof}

Nominal quantities are themselves evaluated with directed intervals.
Their coefficient enclosures give intervals for
$\alpha=\min\{\alpha_{\rm cap},o\exp(-e/2)\}$ and for the monotone
log-transmittance residual. If certified signs satisfy
$R(l)<0<R(h)$ and $h-l\leq2\tau$, the unique median lies in $[l,h]$.
The implementation chooses endpoints within $\tau=10^{-3}$ of its
candidate, so the candidate has that error bound as well.
An uncertified low-precision candidate triggers source-preserving
double reconstruction; failure to certify that result remains unresolved.

Stable evaluation retains amplitude and energy separately. For an uncapped
layer of amplitude $o<1$,
\[
 1-o e^{-e/2}=(1-o)\left[1+\frac{o}{1-o}
       \bigl(-\operatorname{expm1}(-e/2)\bigr)\right].
\]
This permits a small prefix-mass difference to be evaluated with
\texttt{log1p}, even when the effective alpha would round to $1/2$.
Exact binary source offsets are formed before projecting through $U$,
avoiding subtraction of large collinear vectors. Exact source collinearity
gives $e=0$; two layers with effective alpha $1/2$ have the analytic root
$(\rho_1p_1+\rho_2p_2)/(\rho_1+\rho_2)$.

The fixed CPU experiment preserves all 7010 earlier coefficient inputs.
It additionally locks seed 13019 and 92 models before evaluation:
64 ordinary, 16 near-threshold, four exactly critical two-layer,
four tiny-offset two-layer, and four without a finite root.
All 106 finite old/new models receive certified brackets containing the
independent reference, and all four no-root classifications pass.
Seventy-six finite models and four no-root models certify on the low path;
30 finite models require fallback. The maximum observed root error is
$1.361841855\times10^{-5}$ in camera-ray units. This selected test set
does not estimate a real-pixel fallback rate.

An independent world-covariance MP240 check passes 106 endpoint pairs,
4408 used coefficient enclosures, and all four no-root classifications.
An exact rational width audit exposed 81 first-version brackets wider
than $2\times10^{-3}$ by at most $6.680\times10^{-16}$; the archived
first version is retained. Version two rounds the endpoints inward and
passes the unchanged width requirement on the identical inputs.
The cache would contain six F32 factor entries, three F64 means, and two
F32 budgets, nominally 56 bytes before untested device alignment.

These certificates trust experimental \texttt{mpmath.iv} arithmetic
with a 53-bit significand and unbounded exponent~\cite{mpmathiv}.
They are not IEEE FP64 GPU, complete renderer, backward, latency, GT,
or training-quality evidence. Exact source offsets and exponential-tail
signs still require a verified device implementation. Mixed precision
already appears in 3DGS contribution testing~\cite{flicker}; neither
mixed precision nor generic interval bracketing is claimed as new.
The particular median certificate's novelty remains unverified.

\paragraph{A device arithmetic component with a signed tail enclosure.}
The CUDA 12.8 documentation describes its standard mathematical function
error bounds as results of extensive, non-exhaustive tests, rather than
guaranteed bounds~\cite{cudamath}. Thus a guessed ULP expansion of those
functions does not by itself supply a rigorous enclosure. A shared
host/device reference core instead uses basic directed IEEE operations,
range reduction, and explicit series remainders. This uses standard
numerical-analysis ideas, not a claim of new elementary functions.

For $g=\exp(\ell)$ and $b=-\log(1-g)/2$, when $\ell<-40$,
\[
 \log b\in
 [\ell-\log2,\;\ell-\log2+2^{-56}].
\]
Indeed, $57\log2<40$ implies $g<2^{-57}$, while
$1\leq-\log(1-g)/g\leq1/(1-g)$. The additive log correction is bounded
by $-\log(1-g)\leq g/(1-g)<2^{-56}$.
The residual's positive and negative terms can therefore be summed after
a common log scaling, retaining the signs of tails that a binary64
unscaled value cannot represent. Scaled values below $\exp(-700)$ are
enclosed by $[0,2^{-1009}]$, since $1009\log2<700$; they are not deleted.
The required binary64 bounds on $\log2$ are independently checked by
exact rational atanh series and a geometric remainder.

The native Windows host core passes 774 locked scalar checks and all
212 bracket endpoints of the 106 finite models. A separately compiled
CUDA core passes those same 212 endpoints on an RTX 5090 Laptop GPU.
Its five-batch median event time is $222.9916504$ ms, including output
allocation/fill and the endpoint kernel. This is a reference-component
measurement on CPU-provided coefficient intervals and CPU-chosen
brackets; device preprocessing budgets, mass classification, selection,
complete root solving, backward, and full-renderer costs are absent.
The device does not separately test all 774 scalar inputs. These results
support the checked arithmetic component and do not establish rendering
efficiency or GT quality. Compiler and IEEE intrinsic semantics remain
trusted. Sources and all inputs are archived in
\texttt{strict\_interval\_host\_v1} and
\texttt{strict\_interval\_cuda\_v1}.

\paragraph{Device source budgets and independently certified root candidates.}
The subsequent device prototype constructs the interval precision factor
from source means, scales, quaternion or rotation, and camera affine map.
A two-pass modified Gram--Schmidt nominal factor supplies $\widehat Q,U$;
the polar and factor-residual budgets above bound its error without
assuming an exactly unit rounded quaternion or camera rotation.
Both low and high caches enclose all 7322 fixed layer inputs in
58,576 peak/rho/energy/alpha checks. Device terminal-mass classification
passes all 110 models on both paths.

The device then generates its own root candidates and global brackets.
It receives only source parameters and model start/count, with no CPU
candidate, bracket, or MP-root input. Standard mathematical functions
are permitted in candidate generation; every returned finite value
still needs the separate strict residual certificate. Low nominal
coefficients are rounded to F32; a failed certificate triggers a new
source-amplitude/energy F64 candidate. The tolerance is a binary64
number strictly below $10^{-3}$, with endpoints rounded inward.

All 106 finite models pass, with 76 low and 30 high certificates, and all
four no-root classifications are correct. The largest observed root
error is $1.361838888\times10^{-5}$ camera-ray units. A separate MP240
world-covariance inverse check verifies 106 newly device-generated
endpoint pairs, 4408 used coefficient enclosures, exact rational widths,
and the four mass classifications. The fixed seed was already observed
in prior CPU/component development, so this is fixed device development,
not independent unobserved confirmation.

On this diagnostic workload, median CUDA events are $0.441504$ ms for
7322 dual-path source budgets, $705.173584$ ms for the 110-model root
and certificate batch, and $705.568237$ ms combined. Output allocation
and fill are included; host transfers and the full renderer are absent.
This cost has no end-to-end efficiency adoption result. The prototype
does not verify renderer selection, its five-step initial range policy,
world backward, or GT quality. Correct root values do not repair the
unbounded opacity derivatives of the retained thin critical examples.
Cheap sufficient certificates and full rendering costs remain next steps.

\paragraph{Exact reciprocal caching: an isolated local efficiency result.}
A fixed implementation optimization replaces six positive-integer reciprocal
calls in the strict series core by a constant-memory table for $k=1,\ldots,51$.
For each entry, exact rational arithmetic verifies adjacent binary64 bounds
$\ell_k\leq 1/k\leq u_k$ (or equality at a representable reciprocal).
These are the same directed IEEE endpoints used by the original division;
the interval multiplication, all other divisions, series degrees, remainders,
candidate solver, tolerances and inputs remain unchanged.
The separately rebuilt control and variant produce bitwise-identical complete
budget and root arrays, including NaN bit patterns, on the same 7322 layer
inputs and 110 models. The control also matches the archived original outputs.
An independent CPU audit checks the actual compiled table literals, the six
isolated source replacements, archive hashes, output bits, exact widths and
strict residual signs. This inherits the earlier MP240 reference checks
through output identity; it is not a new unobserved-cohort confirmation.

Five alternating paired CUDA-event batches give combined medians of
$705.739014$ ms for control and $385.180573$ ms for the variant, and a
geometric mean paired ratio $0.545837$. The local diagnostic time decreases
by $45.4163\%$ (speed ratio $1.83205$), passing the preset ratio $\leq0.95$
and bitwise-equivalence gates. Output allocation and fill are included;
host transfer, full rendering, backward and training are excluded.
The first generator run failed a substitution-count assertion before CUDA
comparison; its source and failure remain archived, and the correction changes
no gate or input. The successful run is
\texttt{reciprocal\_root\_v2}. Integer reciprocal caching is an implementation
optimization, not an established novel mathematical method. The remaining
strict diagnostic cost does not establish production or GT adoption.

\paragraph{Short sufficient certificates with a retained strict fallback.}
A single fixed short layer keeps the same source budgets, candidate solver,
mass classification, tolerance and input cohort. Its normalized-log expansion
has nine terms and positive remainder at most
$2z^{19}/[19(1-z^2)]$. Its range-reduced exponential uses degree 12 with
absolute remainder at most $3r^{13}/13!$ for radius $r\leq1$.
For radius $r<1/8$, degree-eight log1p and expm1 use the respective bounds
$r^9/[9(1-r)]$ and $r^9/[9!(1-r)]$. Directed basic operations evaluate
all these bounds. Short certificates use the same strict opposite endpoint
signs; failure invokes the original long certificate at the same candidate,
followed, if necessary, by the existing high path. Standard library ULP
assumptions are not used for either certificate.

All 774 saved scalar inputs pass independent Decimal240 and MP240 elementary
checks. All 106 finite models finally use short certificates; 76 use the low
path and 30 use high. There are nevertheless 30 strict calls on failed low
attempts; this is not a zero-fallback result. Four no-root classifications
remain correct. A separate MP240 world inverse check passes all 106 new
endpoint pairs, 4408 used coefficient checks, scaled enclosures and rational
widths. On five alternating paired batches against the reciprocal-table
control, combined CUDA-event medians are $385.132538$ and $245.741669$ ms,
with geometric mean ratio $0.638155$, a local time reduction of $36.1845\%$.
This passes the fixed local gate, not full-render or training adoption.

A separate fixed flat family tests successful strict fallback and unresolved
behavior without altering that cohort or its timing. Eight isotropic single
primitives have peak 1, amplitude $3/4$, and scale $10^j$, $j=0,\ldots,7$;
their exact median is 1. The control certifies four. The hybrid certifies
three through short and one through strict, and retains four unresolved
cases. All returned certificates pass the independent MP240 single-primitive
formula and exact width checks. The family supplies no real-pixel prevalence
estimate or uniform certification guarantee. These are standard series
implementation bounds; academic novelty, scalable full rendering, world
backward and GT quality remain unverified.

\paragraph{Actual complete-image contributor capacity.}
Before moving the local solver into the renderer, an isolated counter is
inserted at the actual RGB acceptance and first median-iteration acceptance
sites of the original affine CUDA forward. The same-source control and
instrumented versions preserve all original RGB, alpha, normal, depth,
slope and radii bits on all 24 fixed development views, comprising
10,834,488 pixels. Counting has no list truncation or sampling. Its generated
source changes only integer counters and additional output channels.
All initially eligible pixels have equal RGB and first-median acceptance
counts; original invalid depth values remain unchanged.

There are 9,165,957 initially eligible and 9,145,921 finally valid pixels.
Accepted-list lengths have median 88, 90th percentile 155, 99th percentile
234 and maximum 461 in both populations. A capacity of 192 misses 334,083
initially eligible pixels and 334,072 finally valid pixels, or 3.652688\%
of the latter. Capacity 256 still misses 41,536 initial pixels. Although 512
covers this observed cohort, it supplies no universal capacity guarantee.
The local 192-layer solver therefore cannot be used as a complete-image
implementation. Dynamic lists must preserve every accepted contributor.

A separate CPU audit recomputes the saved complete depth/slope bit checks,
counts and histograms and checks 669 sources, 50 artifacts, both generated
source trees, binaries and vendor identities. RGB/alpha/normal/radii are
fully compared at runtime but not retained for a second offline bit audit.
Five alternating native-render pairs per view give a geometric mean
counting-overhead ratio of 1.025855. This measures the counter, not certified
rendering or training. Native forward excludes appearance postprocessing;
all views were already observed in development. The initial UTF8-source
audit failed under the Windows default GBK decoder; an explicit UTF8 read
repairs that verifier without changing any experiment data or gate.

\paragraph{Dynamic lists preserve the complete accepted topology.}
A subsequent private native-forward diagnostic first reads the frozen
acceptance counts and allocates int64 prefix offsets and int32 ID/position
pairs. The original first-iteration acceptance site writes each actual
primitive ID and its position in the original tile point-list. No fixed
192/256/512 list cap is used. Each write checks both its pixel span and
the total allocation; an overrun attempt sets bit 1, and an actual/expected
count mismatch sets bit 2. Nonzero status is a failure, not a certificate.

All 24 views retain 870,178,922 associations, exactly matching the prior
complete initial-count total. Original native image outputs and counters
remain bitwise identical. Every slot is filled and every write/checker
status is zero. An independent GPU checker verifies increasing positions,
the actual contributor prefix and tile bounds, and the matching point-list
ID. A separate CPU audit rechecks every saved association with bounded
memory, including all offsets, counts, ID ranges, within-pixel ordering and
point-list provenance, with no sampling. It also checks 36 sources, 54
artifacts, 18 generated sources, binaries, vendor headers and the archive.
Complete median/counter arrays are audited offline; the original runtime
RGB/alpha/normal/radii comparisons are not a second saved-array bit audit.

Two separately stored wrong-allocation probes remove or add one slot at
the first nonzero pixel. They respectively return status 3 and status 2
only at that pixel; over-allocation leaves exactly one sentinel, and all
original image bits remain unchanged. These failed inputs do not count
as successful captured lists.

Per-view records occupy 157--383 MB. Five alternating capture/count pairs
with allocated lists reused give a geometric mean ratio 1.074844.
A separate complete workflow, including both forward passes, scan,
allocation and checker, takes 23.680--50.505 ms in one CUDA-event observation
per view. Host gaps and synchronization are included; transfer/storage is
separate. This is a correct topology interface, not a rendering speedup,
certified full-image root solver, backward or GT result. Shared source
parameter caches and ray-parallel certification remain to be implemented.

\paragraph{Shared query-independent geometry budgets.}
The frozen preprocessing is next split at the query boundary: source
inverse, nominal QR, polar/residual errors and mean budgets are cached per
primitive and camera; the original query-dependent statements are retained
verbatim. The diagnostic cache has 28 doubles (224 bytes), including both
statuses, both upward-F32 factor budgets, the upward-F32 mean budget,
nominal and interval means, both six-entry factors and padding. It is not
the proposed 56-byte production layout. True mean intervals are retained
for the original proven-zero-energy alignment branch.

On all 7322 source rows, complete two-mode budget arrays and all 106 finite
plus four no-root outputs remain bitwise identical to the frozen reference,
including NaN payloads. Eight repeated query sets (58,576 rows) also match.
Five alternating original-repeat versus build-once/cache-decode pairs give
geometric mean CUDA-event ratio 0.830169, a local time reduction of
16.9831\%, passing the fixed bitwise and timing gates. Allocation/fill and
the query ID-range checks are included; host input transfer and full
rendering/training are absent. Earlier MP240 evidence is inherited through
output identity, not rerun as a new unobserved cohort.

All primitives in each of the 24 real development cameras also receive
finite, status-zero caches. Source/model/camera identities match the prior
CSR inputs; no primitive is dropped by visibility. Per-view cache memory
is approximately 114--236 MB, with a separate 288-byte source row per
primitive in this diagnostic. The real cache hashes, error-budget
statistics and five build event/wall observations are recorded at runtime;
the real arrays are not retained for a second offline full MP audit.
The independent audit checks the saved local arrays, unchanged query
body, 25 sources, 28 artifacts, 12 generated files and real build metadata.
This does not yet decode all real CSR associations or certify complete
images, world backward, GT or training performance.

\paragraph{Actual complete CSR coefficient decoding.}
The next private forward adds only two output channels for the actual
kernel pixnf coordinates, retaining all original image and counter bits.
The amplitude is read from the actual conic-opacity buffer, not assumed
equal to the input opacity. The query bits match all nine earlier anomalous
ray captures. Source/model/camera identities and every accepted-count map
match the saved CSR topology.

All 870,178,922 associations in the 24 views are decoded in batches of
65,536 using int64 offset upper-bound lookup and the frozen shared-cache
query body. Both modes return status zero and finite valid outputs for
every association. The maximum temporary double output is 27,262,976 bytes;
cache, CSR and query/amplitude residency are additional. Batch coverage,
status counts, output hashes and event/wall costs are saved, but the full
approximately 362 GB double streams are not retained to recompute every
hash offline.

Predeclared 48 evenly spaced records per view plus all earlier anomalous-ray
associations yield 1924 probes. Complete probe arrays are bitwise identical
to direct preprocessing. Independent MP80/140 world-covariance inverse
calculations agree and pass 15,392 two-mode peak/rho/uncapped-energy/capped-
alpha interval checks. The audit checks 39 sources, 74 artifacts, generated
query/strict sources, binaries and archives. These fixed probes do not
constitute individual MP checks of all 870 million records or image root
certificates.

The diagnostic dual-mode decode costs 597.529--2164.430 ms per view in
summed CUDA events, before CPU transfer, hashing and statistics. This is
too costly for direct method adoption. In each mode, 21,577 nominal
double amplitude-times-exponential calculations underflow to zero while
the strict source intervals remain valid. These records are retained;
nominal zero does not prove mathematical zero or valid gradients. The
measurement precedes conversion of candidate alpha to F32. Cheap sufficient
tests, selective strict fallback and ray-parallel roots remain next steps;
world backward, GT and training gains are unverified.

\paragraph{Direct F32 product sufficient conditions.}
The next isolated diagnostic retains the low-mode directed double geometry
and replaces the alpha exponential and endpoint transmittance product with
outward F32 intervals. A warp processes each complete CSR ray; no 192-node
cap or contributor pruning is introduced. For $x\leq0$, range reduction
$x=k\log2+r$, with $|r|\leq1$, uses the degree-eight exponential polynomial
and the absolute remainder $3|r|^9/9!$, since $e^{|r|}<3$. The adjacent F32
enclosure of $\log2$ and the reciprocal coefficient enclosures are verified
using exact rational arithmetic. For $x<-80$, the enclosure is
$[0,2^{-115}]$, justified by $115(\log2)_{\rm upper}<80$;
underflow is not identified with an exact mathematical zero. Multiplication,
division and square root use CUDA directed intrinsics with FTZ disabled;
no empirical standard-exponential ULP assumption enters the certificate.
An uncertain peak-side decision uses the hull of both branch intervals.

\begin{proposition}[Including public-depth rounding in the certificate]
Let $z$ be the reported finite positive F32 camera depth, promoted exactly,
and let $q=(q_x,q_y,1)$ contain the actual F32 kernel query coordinates.
Suppose directed arithmetic encloses $c=z\|q\|$ in $[c_-,c_+]$.
Choose $0<\tau<10^{-3}$ and form
\[
 \ell=\operatorname{round}_{+\infty}(c_+-\tau),\qquad
 u=\operatorname{round}_{-\infty}(c_-+\tau).
\]
For the fixed physical Gaussian contributor list, if $\ell<u$ and valid
transmittance enclosures satisfy
$T(\ell)_{\rm lower}>1/2>T(u)_{\rm upper}$, a unique median root $t^*$ exists
and $|t^*-c|<\tau$.
\end{proposition}
\begin{proof}
The strict inequalities imply that at least one layer has positive
effective amplitude. Each valid positive-amplitude Gaussian layer is
strictly decreasing in transmittance over every nonempty interval, even
if its derivative vanishes at its peak. The full product is continuous and
strictly decreasing, hence has a unique root between $\ell$ and $u$.
Directed rounding gives $\ell\geq c_+-\tau\geq c-\tau$ and
$u\leq c_-+\tau\leq c+\tau$, proving the claimed strict error bound.
\end{proof}
The implementation uses the previously rationally verified tolerance
strictly below $1/1000$. Invalid original depths are separate ``no original
candidate'' states, not mathematical no-root classifications; a failed
sufficient test remains unresolved.

All 10,834,488 pixels and their 870,178,922 associations are accounted for
in 24 frozen development views. Of 9,145,921 valid original depths,
8,583,586 (93.8515\%) are certified and 562,335 remain unresolved.
All rays retain their original contributors; no bad-input flag is observed.
Complete compact outputs are saved. The independent audit checks 25 sources,
76 artifacts and 4104 scalar enclosures. Predeclared 64 evenly spaced
eligible pixels per view plus the nine prior anomalies yield 1545 complete
rays, 145,859 source associations and 583,436 coefficient checks against
independent MP80/140 world-covariance inverse calculations. All sampled
coefficient and endpoint bounds pass, including 1410 certified probe rays.
This is not individual MP verification of all 870 million associations.

An explicitly exploratory follow-up analyzes all 132 unresolved rays within
that fixed probe set. At the exact public-depth center plus/minus $1/1000$,
independent MP80/140 signs locate 18 roots before and 114 roots after the
error band; none is an interval-width-only failure in this subset. All nine
previous anomalous rays are included. This conditional subset cannot
estimate the population error rate or reconstruction quality.

The median of three repeated diagnostic events per view ranges from
57.714 to 134.755 ms, including compact output allocation/fill and kernel.
Cache construction and existing CSR/query residency are additional.
The diagnostic is substantially cheaper than the preceding full double
output audit but remains slower than the original renderer; it is not an
end-to-end speed result. A dynamic strict root correction for unresolved
pixels, world VJP, full GT comparisons, training and novelty remain open.
Reproducible sources and results are recorded in
\texttt{FLOAT\_PRODUCT\_PROTOCOL.md} and
\texttt{float\_product\_audit\_v1.json}.

\paragraph{Dynamic correction of all unresolved original-valid rays.}
The next isolated implementation decodes all 562,335 previously unresolved
original-valid rays, comprising 28,196,963 source associations. Dynamic
int64 segments retain all contributors: 320 rays exceed the previous
192-layer solver cap, and the selected maximum is 372. A warp uses the
stable nominal source-amplitude/energy prefix shift and scaled tails to
propose a double root. A rounded numerical zero can only propose a value.
The public output is then rounded to F32 camera depth and certified around
its real query-norm center by the preceding proposition. F32 products,
short bounded scaled-log series and the strict bounded series are tried
in sequence; a failed low geometry mode is redecoded in the high mode.
The search uses global peak/reach brackets, rather than claiming identity
with the legacy initial numerical search interval. The fixed contributor
list and physical target transmittance are preserved.

The full selected set yields 562,079 certified corrected depths, 256
certified no-finite-root classifications, and zero remaining unresolved
states. Of the finite results, 562,050 use low geometry and 29 use high
geometry; 562,053 finish with F32 product certificates and 26 with short
scaled-log certificates. Strict fallback is executed, even though no final
accepted result requires that last certificate. Every certified public
depth is positive. The original 8,583,586 certified values are not touched;
the 1,688,567 originally ineligible/no-candidate pixels are not filled.
Thus this does not claim a valid median for every image pixel.

An independent audit checks 30 sources, 74 artifacts, generated headers,
actual binaries, source archives and full output/CSR accounting. The
predeclared union of the earlier 132 unresolved probes and 32 evenly
spaced unresolved pixels per view contains 891 complete rays and 51,524
source associations. Independent MP80/140 world-covariance inverses pass
206,096 coefficient bounds; MP240 stable scaled-log arithmetic verifies
the saved public F32 endpoint signs and depth bands. All nine earlier
anomalous rays are corrected and certified. This fixed probe set contains
no no-root case, so its success alone does not validate that rare branch.

A separately declared exploratory conditional audit therefore checks the
entire union of all 256 no-root classifications, all 29 high-mode finite
results and all 26 short-certificate results: 299 rays, 5920 associations
and 23,680 coefficient checks. All pass, including the finite public
endpoint tests. Every no-root case has true final transmittance above
$1/2$, with minimum excess approximately $4.51860\times10^{-5}$.
This is existence in the fixed physical Gaussian model, not absence of
a real-world surface. A separate source/artifact audit checks the entire
conditional selection, query/amplitude/CSR mappings and unchanged saved
root bits. Conditional rare-branch checks are not random holdout evidence.

The per-view sum of low/high root-event medians over three observations
ranges from 5.822 to 204.906 ms; single decode-event sums range from
1.625 to 3258.099 ms. The first-view 3258.099 ms observation is retained
without attributing it to an unverified loading hypothesis. These selected
ray diagnostics omit additional cache construction, CPU compaction,
original-candidate certification and contributor capture/full-renderer
costs, and are not an end-to-end speed result. Added peak GPU allocation
is about 774 MB, on top of existing residency. Source world VJP, full GT
reconstruction, training quality and novelty remain unverified.

\paragraph{World parameter chain and the point used by the adjoint.}
For a fixed camera ray $o+td$, let $v=\mu-o$, $P=\Sigma^{-1}$,
$a=d^\top Pd$, $p=d^\top Pv/a$, $\rho=\sqrt a$, and $w=v-pd$.
With the ray held fixed, exact matrix calculus gives
\[
 \delta p=\frac{d^\top P\delta\mu+d^\top\delta P\,w}{a},\quad
 \delta\rho=\frac{d^\top\delta P\,d}{2\rho},\quad
 \delta E=2(Pw)^\top\delta\mu+w^\top\delta P\,w.
\]
The last identity uses $d^\top Pw=0$, so the variation of the projected
peak does not enter the energy variation separately. If layer partials
are denoted by $h_p,h_\rho,h_E$, their precision upstream is
\[
 G_P=\frac{h_p}{2a}(dw^\top+wd^\top)
       +\frac{h_\rho}{2\rho}dd^\top+h_Eww^\top.
\]
For the actual invertible quaternion-polynomial matrix $R$, not an assumed
exact orthogonal rotation, write $F=S^{-1}R^{-1}$ and $P=F^\top F$.
The corresponding parameter partials are
\[
 h_{\log s_j}=-2f_j^\top G_P f_j,\qquad
 G_R=-2PG_PR^{-\top},
\]
where $f_j^\top$ is row $j$ of $F$. These standard chain identities are
bookkeeping for the implementation, not a novelty claim. The root VJP is
$-H_\theta/H_t$, and its public camera-depth counterpart additionally
divides by the fixed actual query norm.

The next CPU audit uses all earlier 891 corrected-ray probes and 51,524
source nodes. It defines a continuous local extension through the actual
captured F32 mean, filtered scales, quaternion and amplitude. Filtered
scales use the ratio of the raw exponential/filter activation to its base
value; quaternion values add the change in normalized raw quaternion;
amplitudes multiply the sigmoid, 3D-filter and screen-compensation base
ratios. Thus the base physical curve is retained while position, raw
log-scale, quaternion and opacity-logit chains are differentiated. This
is neither a derivative of the literal discrete IEEE program nor a
silent replacement by the earlier unrounded raw-physical model.
The contributor set, filter, camera and query are held fixed.

MP140 world-covariance inverses and analytic parameter chains, followed
by MP240 stable residual/IFT arithmetic, supply an independent reference.
The Torch implementation receives only existing certified endpoints and
source parameters; it polishes its own roots. All 891 references are
regular and all stored gradient arrays are finite. The predeclared
relative-L2 $10^{-5}$ and per-block $10^{-4}$ gates pass for 858 rays when
the adjoint is evaluated at the existing GPU internal candidate, for 345
at the public F32 depth center, and for all 891 after further root polishing
in either inverse-factor or dense-inverse arithmetic. The maximum complete
ray relative-L2 errors are respectively $4.47157\times10^{-3}$,
$4.48923\times10^{-3}$, $1.86370\times10^{-7}$ and
$9.86379\times10^{-7}$. Both polished branches pass; no general superiority
or training improvement follows. Public-center errors include the earlier
candidate error as well as output rounding.

An independent source/array audit checks 13 sources, 51 artifacts,
all base/root mappings and every saved per-ray/block statistic. All actual
views in this batch have screen kernel zero, so they do not empirically
validate a nonzero screen-compensation derivative. A general-helper V2
fixes the frozen compensation baseline for gradient-disabled perturbed
forward calls; all 891 real polished gradients remain bitwise unchanged.
Three explicitly synthetic kernel-zero/nonzero cases then pass complete
root central differences, with final relative errors below
$5.74\times10^{-10}$. These examples are not real reconstruction evidence.
Retaining and refining an internal root is therefore supported by this
development audit, but full CUDA backward, topology-changing derivatives,
GT reconstruction and training quality remain open. A separate fixed-model
49-view reconstruction protocol is recorded before its results exist.

\paragraph{Full frozen-model reconstruction: quality gate fails.}
The preregistered follow-up now evaluates the original scan37/55/63
checkpoints on all 49 training views each, rather than only the earlier
24 development views. All 147 original RGB/alpha/normal/radii/depth/slope
outputs match bitwise. The workflow accounts for 66,361,239 pixels and
5,346,733,664 accepted associations. Among 56,594,143 originally valid
pixels, 52,932,318 retain their sufficient certificates; the remaining
3,661,825 yield 3,659,782 certified corrections, 2043 mathematical no-root
classifications and zero unresolved results. Both branches retain the
same 24,152,107 valid pixels after the identical ground-truth mask.
Thus the no-root invalidations do not reduce masked coverage in this test.

The branches use identical RGB, masks, camera intrinsics/extrinsics,
view order, voxel size $0.002$, depth scale $1$, depth maximum $8$,
truncation multiplier $4$, block count $50,000$ and largest-component
postprocessing. Names 0000--0048 are mapped to their actual official
camera indices; both meshes share one camera-center alignment transform.
Three fixed official evaluator seeds per branch and scene give 18 runs.
They are evaluation-sampling repeats, not independent training repeats
or a statistical confidence statement.

The three-seed mean overall errors, in official DTU millimetres, are
\begin{center}
\begin{tabular}{lrr}
\toprule
Scene & Original & Certified correction\\
\midrule
scan37 & 0.4605930275 & 0.4605502386\\
scan55 & 0.3139134531 & 0.3138697150\\
scan63 & 0.7369190101 & 0.7369960982\\
\bottomrule
\end{tabular}
\end{center}
The ratio of three-scene mean errors is $0.9999937549$, approximately
$0.0006245\%$ lower, which fails the preregistered $1\%$ improvement
threshold. The scene-wise mean surface-to-data error differences are
$+0.000023823$, $+0.000033254$ and $+0.000505426$ mm; all also fail the
nonpositive-difference gate. These small changes are not claimed as
significant reconstruction improvements.

An independent audit checks 97 sources, actual binaries, every saved
image/root/mask and certificate-accounting array, six mesh outputs and
all 18 evaluation sources/results. This accounting is not individual MP
verification of all 5.35 billion associations. Scene wall costs are
approximately 516, 299 and 370 seconds including storage, CPU TSDF,
extraction and evaluation waits; complete per-stage costs are retained.
The numerical correction remains a diagnostic reference, and is rejected
as a demonstrated frozen-model quality optimization. Source/model files
and earlier results remain unchanged. A subsequent candidate objective
must be assessed separately, not credited with results from this failure.

\section{Vacancy and transmittance are conditionally equal}
\begin{proposition}[Exact history gap]
For every finite $t$,
\[
 \Delta(t):=\log\frac{V(t)}{T(t)}
 =\sum_{i:p_i<t}\log\frac{1-g_i(t)}{1-\alpha_i}\geq0.
\]
Under the strict amplitude assumptions, equality holds if and only if no peak
has been passed. Consequently, $T(t^*)=q$ implies
$V(t^*)=q\exp(\Delta(t^*))$, which need not equal $q$.
\end{proposition}
\begin{proof}
Before a peak, $v_i/T_i=1$. After it,
$v_i/T_i=(1-g_i)/(1-\alpha_i)>1$ because $g_i(t)<\alpha_i$.
Taking the logarithm of the product proves the identity and equality criterion.
\end{proof}
This makes explicit the near-side condition in the GGGS depth discussion.
It does not contradict the equivalence within that condition.
For example, two primitives with $\alpha_i=0.4$, peaks $0,1$, and widths
$0.03$ have final transmittance $0.36$, hence a median root. Yet
$V(t)>0.5$ everywhere: indeed each $v_i(t)\geq\sqrt{0.6}$, so
$V(t)\geq0.6>0.5$. A more precise numerical evaluation is unnecessary for
this counterexample. Reversing a ray can therefore change its median without
producing a common $V=0.5$ surface.
Indeed symmetry gives $T(1/2)=0.6$ exactly, so the forward median lies after
the midpoint, while the reversed world-space median lies before it.

\subsection{A squared history gap removes the median singularity}
Fix the number of layers and use independent coefficient coordinates
$\theta=(\alpha_i,p_i,\rho_i)_{i=1}^N$, with $\rho_i=1/\sigma_i$.
Consider the open domain
\[
 \mathcal D=\{0<\alpha_i<1,\ \rho_i>0,\
 \Omega:=-\textstyle\sum_i\log(1-\alpha_i)>\log2\}.
\]
Write $t_*(\theta)$ for the unique finite median,
$\ell=\log V$, $D(\theta)=\ell(t_*,\theta)+\log2$, and $L=D^2$.
The preceding identity gives $0\leq D<\log2$.

\begin{theorem}[Continuous first derivatives despite a singular median]
On $\mathcal D$, $D$ is locally Lipschitz and $L$ is continuously
differentiable. The singular median set is exactly
\[
 \mathcal S=\{p_1=\cdots=p_N,\ \Omega=\log4\}.
\]
On $\mathcal S$, $D=L=0$ and $\nabla L=0$. Outside $\mathcal S$,
\[
 \nabla D=\ell_\theta-rH_\theta,\qquad
 r=\frac{\ell_t}{H_t}\in[-1,1],\qquad
 \nabla L=2D(\ell_\theta-rH_\theta),
\]
where partial derivatives on the right hold $t$ fixed at $t_*$.
No global bound is asserted as $\alpha_i\to1$ or $\rho_i\to\infty$.
\end{theorem}
\begin{proof}
Each layer attenuation is continuously differentiable, strictly increasing
apart from a zero derivative at its peak, with limits zero and
$-\log(1-\alpha_i)$. Hence $H(t,\theta)=\log2$ has a unique finite root.
Given any parameter point, choose two finite strict bracketing endpoints.
Joint continuity preserves their strict signs nearby. Applying the same
argument at endpoints arbitrarily close to the unique root proves that
$t_*$ is continuous.

Let $x_i=(t-p_i)\rho_i$, $u_i=\alpha_i e^{-x_i^2/2}$ and
$b_i=u_i/[2(1-u_i)]>0$. Direct differentiation yields
\[
 \ell_t=\sum_i b_i x_i\rho_i,\qquad
 H_t=\sum_i b_i|x_i|\rho_i.
\]
Thus $H_t=0$ if and only if all peaks coincide with $t$.
At that point $H=\Omega/2$; the median equation gives
$\Omega=\log4$. Conversely these conditions give a singular median.
Also $\ell=-\Omega/2=-\log2$ there, so $D=0$.
At a regular median, the implicit function theorem gives
$\nabla t_*=-H_\theta/H_t$ and the stated formula. The triangle
inequality gives $|r|\leq1$.

Choose a closed Euclidean ball contained in $\mathcal D$. Continuity
bounds $t_*$ on this ball. All fixed-$t$ first derivatives of $H$ and
$\ell$ are continuous, including where $t=p_i$: the two layer branches
have the same value and first derivatives there. On the compact set of
these parameters and bounded root depths, these derivatives are uniformly
bounded. Therefore $\|\nabla D\|$ has a uniform bound away from
$\mathcal S$ within the ball, independently of how small $H_t$ becomes.

For completeness this bound extends across $\mathcal S$, rather than
merely holding on its complement. A line segment in the ball either has
at most one point with all peaks equal, or has equal peaks throughout.
In the latter case, along its affine opacity coordinates,
\[
 \frac{d^2\Omega}{ds^2}
 =\sum_i\frac{(\alpha_i')^2}{(1-\alpha_i(s))^2}.
\]
If any opacity direction is nonzero, strict convexity gives at most two
intersections with $\Omega=\log4$. Otherwise the segment either misses
$\mathcal S$ or lies entirely in it. Integration of the bounded regular
derivative along the finitely many complementary subsegments, and
continuity at their endpoints, bounds the change of $D$ by a constant
times segment length. If the segment lies in $\mathcal S$, $D$ is zero.
This proves local Lipschitz continuity.

At $\theta_s\in\mathcal S$, this bound gives
$L(\theta_s+v)=O(\|v\|^2)$, so $L$ is differentiable there with
zero derivative. On approaching $\mathcal S$ from its complement,
$\nabla L=2D\nabla D\to0$, by continuity of $D$ and the same local
bound. Away from $\mathcal S$, all terms are continuous with $H_t>0$.
This proves continuous differentiability everywhere in $\mathcal D$.
\end{proof}

For coefficient implementation let $s_i=1$ before the peak and $s_i=-1$
after it, and $c_i=b_i(1+rs_i)$. At a regular root,
\[
 D_{\alpha_i}=-\frac{c_i}{\alpha_i}
       -\frac{r\ind[t_*>p_i]}{1-\alpha_i},\quad
 D_{p_i}=-c_i x_i\rho_i,\quad
 D_{\rho_i}=\frac{c_i x_i^2}{\rho_i}.
\]
Normalize the log densities to form $r$; separately computing a large
$1/H_t$ is unnecessary. If $A$ and $B$ are the respective sums of
after-peak and before-peak nonnegative densities, compute
$1+rs_i$ as $2A/(A+B)$ before a peak and $2B/(A+B)$ after it.
This avoids cancellation when $r$ rounds to $1$ or $-1$. Also evaluate
$b_i/\alpha_i=e^{-x_i^2/2}/[2(1-u_i)]$ without division by a tiny
$\alpha_i$. Real source cases include positive effective amplitudes
which round to zero in Float64. The audited CPU implementation therefore
retains a separate log-amplitude input for density normalization; this
is not an amplitude-only interface guarantee.

\begin{remark}[Smoothness does not establish useful optimization]
When all peaks coincide, the exact identity is
$D=[\log4-\Omega]_+$. Thus $L$ is a squared hinge and generally not
twice differentiable: for one layer at $\alpha=3/4$ its second opacity
derivatives from below and above are $32$ and $0$.
More adversely, take two layers with $\alpha_1=\alpha_2=1/2$,
$p_1=-1$, $p_2=1$, and $\rho_1=\rho_2=\rho$.
Their exact median is zero, $r=0$, and
\[
 D=\log(2-e^{-\rho^2/2}),\qquad
 L\longrightarrow(\log2)^2,
 \qquad \|\nabla_\theta L\|=O(\rho^2e^{-\rho^2/2})\longrightarrow0.
\]
The loss can therefore remain large with exponentially vanishing
gradients. It is not a guaranteed remedy for separated surfaces.
This theorem concerns continuous independent layer coefficients; cap,
visibility changes and the discrete floating-point input program are
outside its differentiability claim. Neither novelty nor improved
training or reconstruction follows from the theorem.
\end{remark}

The fixed numerical audit passes all eight symbolic identities and
16 regular synthetic cases: whole-root-composed finite differences
at steps $10^{-30}$ and $10^{-35}$ have maximum relative error
$5.41\times10^{-47}$. Four adverse thin-gap cases use 420-digit
arithmetic and tail-dependent steps down to $10^{-147}$, with two
step sizes; they confirm, rather than remove, the vanishing-gradient
counterexample. At $\rho=20$, $L\simeq0.480453$ while the coefficient
gradient norm is $2.72\times10^{-85}$.

The real-source CPU check retains all 891 fixed development rays,
51,524 nodes and nine previous anomalies. It rederives coefficients
by a 140-digit world covariance inverse, refines the reference median
at 240 digits, and checks the independent history identity. All 891
Float64 coefficient gradients pass both the whole-ray $10^{-5}$ and
per-block $10^{-4}$ relative L2 gates; maximum errors are respectively
$5.70\times10^{-12}$ and $1.16\times10^{-11}$. Five positive
amplitudes round to zero in Float64. In contrast, evaluation at the
existing public Float32 depth center passes only 264 of 891 rays,
with maximum whole-ray relative error $0.004158$.
An initial implementation stopped on a zero-rounded-amplitude NaN;
its corrected successor was finite but failed 108 block gates through
$1-r$ cancellation. Both failures and their original sources remain
saved. The final artifact audit checks all stored arrays, 55 sources
and 53 artifacts. These checks establish neither a population error
rate nor a complete raw-parameter/CUDA training derivative or an
improvement over ground truth. The squared history gap remains a
mathematical supplement and a candidate objective with a serious
optimization limitation.

\subsection{A continuation when no finite median exists}
The coefficient domain can be enlarged without altering any finite-root
loss. Define, on $0<\alpha_i<1$ and $\rho_i>0$ with fixed finite $N$,
\[
 \overline L(\theta)=
 \begin{cases}
  L(\theta),&\Omega>\log2,\\
  (\log4-\Omega)^2,&\Omega\leq\log2.
 \end{cases}
\]
The second branch penalizes insufficient attenuation; it does not assign
zero loss to a missing root.
\begin{theorem}[Continuous first derivatives through the mass boundary]
The extended loss $\overline L$ is continuously differentiable on this
entire open coefficient domain. It agrees exactly with the original
loss wherever a finite median exists. On the insufficient-mass branch,
\[
 \overline L_{\alpha_i}=-\frac{2(\log4-\Omega)}{1-\alpha_i},
 \qquad \overline L_{p_i}=\overline L_{\rho_i}=0.
\]
\end{theorem}
\begin{proof}
Only $\Omega=\log2$ needs an additional argument. Fix any parameter
point on that boundary and put $p_{\max}=\max_i p_i$. At every finite
depth, each positive Gaussian layer has strictly less attenuation than
its asymptotic total. Consequently
$H(p_{\max},\theta)<\Omega=\log2$.
The maximum of finitely many continuous peak coordinates is continuous,
so this strict inequality persists in an open neighborhood of the point.
Whenever a finite median exists in this neighborhood, monotonicity places
it strictly beyond every peak. On that part of the neighborhood,
$H=\Omega+\ell$, whence the median equation gives
$D=\ell+\log2=\log4-\Omega$. Thus on both sides of the boundary
$\overline L$ equals the same smooth formula on an entire neighborhood,
not merely in a limiting-value sense. Below the boundary that formula is
smooth; above it the preceding theorem applies. Differentiating the
smooth formula proves the derivative statement.
\end{proof}

The extension audit checks 16 fixed random models on both sides of the
mass boundary against independent root-composed or product-mass finite
differences at 240 digits; maximum relative error is
$3.34\times10^{-56}$. Sixteen two-sided seam cases and the exact
single-layer $\alpha=1/2$ boundary agree with the tail-region formula.
The two previously missing full/geometry updates are the same real ray,
with independently computed $\log2-\Omega$ approximately $0.0025415$
and $0.0025419$. The extended losses are respectively $0.48398274$
and $0.48398332$, with Float64 value discrepancies below
$6.0\times10^{-15}$. Their raw directional derivatives agree with
independent 240-digit composed differences within $2.4\times10^{-14}$
relative error. These two updates validate this diagnostic extension;
their original undefined branch means are not overwritten, and no
full-branch MP or training benefit is inferred.

For fixed geometry, if an earliest peak has effective amplitude at least
$3/4$, its own attenuation at that peak is at least $\log2$.
The median therefore lies no later than any peak and $L=0$.
This sufficient condition needs no coincident peaks. It shows that a
zero history gap alone does not identify aligned geometry. Whether a
raw opacity can attain that effective amplitude depends on transverse
energy, filtering and the amplitude cap; this coefficient-level
observation is not a claim about every source ray.

\subsection{Raw-parameter and shared-primitive diagnostics}
The subsequent fixed-source CPU audit includes all 891 rays and 51,524
nodes. It chains the coefficient adjoint through positions, filtered
log-scales, normalized raw quaternions, opacity logits and active
amplitude caps. Source cameras, filters, contributor IDs and visibility
are fixed, using the previously documented continuous extension through
captured inputs. The factor implementation passes all 891 comparisons
against an independent 140/240-digit world-inverse analytic chain:
maximum whole-ray relative L2 is $4.54\times10^{-11}$, and maximum
parameter-block error is $2.54\times10^{-10}$. A dense Float64 covariance
implementation passes only 890 rays, with maximum whole-ray error
$2.61\times10^{-5}$; the original combined gate remains failed.

Of 42 prespecified re-solved directional checks on 21 rays, 35 pass at
the prescribed final step $10^{-5}$. A separate follow-up retains all
seven failures and uses an independent 240-digit composed raw forward:
all seven agree at micro steps, with last relative errors below
$6.0\times10^{-17}$. These selected diagnostic checks do not replace
the original failed finite-step verdict. Three additional nonzero-screen
kernel cases are synthetic and kept separate. In this real development
set, 561 rays have loss at least $0.25$, none of which meet the
prespecified weak scaled-gradient threshold $10^{-8}$. Three zero-loss
rays have zero gradient. This does not remove the thin-gap counterexample
or estimate its population frequency.

A common update then links actual primitive IDs across the rays:
24,007 unique primitives, of which 6,278 repeat, with maximum
multiplicity 137. Identical raw inputs and scaling for repeated IDs are
verified; zero-step losses match bit for bit. At equal normalized
coordinate step $10^{-3}$, the mean loss changes from $0.2891820143$
to $0.2889393258$ for the full direction and $0.2889393275$ for the
geometry-only direction. The opacity-only direction gives
$0.2891811147$. These are fixed-ray loss changes, not ground-truth
geometry gains. Under the declared parameter scaling, about $99.14\%$
of aggregate squared gradient energy is in quaternion coordinates;
this is not an optimizer-independent attribution of training effects.
At step $10^{-2}$, both full and geometry-only updates lose one finite
root. Their original branch means are left undefined, preserving all
891 rays in the accounting.

The original shared central-difference gate has relative error $5.35\%$
at step $10^{-5}$, dominated by one retained old anomaly. Follow-up
evaluation of all 891 rays at $10^{-9}$ yields relative error
$8.58\times10^{-6}$, whereas $10^{-11}$ worsens to
$2.23\times10^{-4}$. At the dominant point, independent 240-digit
forward differences give relative errors approximately $5.00\times10^{-12}$
and $4.67\times10^{-16}$ at these latter two steps. Thus both truncation
and Float64 resolution limit the finite-difference diagnostic; the
original gate is not relabeled as passed. None of these local updates
establishes CUDA correctness, full-visibility training behavior,
photometric or reconstruction improvements, or novelty.

\subsection{A GPU implementation of the extended loss}
Loss evaluation can avoid solving two classes of roots. Monotonicity
and the exact history identity imply
\[
 H(p_{\min})\geq\log2\ \Longrightarrow\ \overline L=0,
 \qquad
 H(p_{\max})\leq\log2\ \Longrightarrow\
 \overline L=(\log4-\Omega)^2.
\]
The second implication includes the insufficient-mass domain; when
the median exists it is beyond all peaks, giving the same formula.
Otherwise the median lies between the extreme peaks. The diagnostic
CUDA implementation evaluates these branches, independently bisects
the interior interval, and returns coefficient seeds for
$(p,\rho,E,A)$; a Float64 factor graph chains these through the full
captured-source raw model. It receives neither MP roots nor CPU adjoints.
Separate before/after log-density normalization forms the bounded ratio
without first computing a large median Jacobian. Compensated prefix
products and sums stabilize the scaled signed residual. These ordinary
Float64 branch signs are not interval certificates.

One numerical correction uses a standard linear-algebra identity.
For invertible $F=S^{-1}R^{-1}$, world displacement $v$ and direction $d$,
\[
 E=\frac{\|Fv\mathbin{\times}Fd\|^2}{\|Fd\|^2}
  =\frac{\|SR^T(v\mathbin{\times}d)\|^2}
   {(\det S\det R)^2\|Fd\|^2}.
\]
This follows from
$(Fx)\times(Fy)=\det(F)F^{-T}(x\times y)$ and
$F^{-T}=SR^T$; it does not assume that captured $R$ is orthogonal.
In the aligned synthetic fixtures this form retains an exact zero
transverse vector, while subtracting a projected residual introduces
tiny nonzero derivatives. Screen compensation and its derivatives are
also identically zero for kernel zero, and are evaluated by that exact
branch. Neither identity is claimed as new.

The first GPU version passes all 891 real development rays, but only
28 of 32 subsequent synthetic raw models pass the unchanged per-block
rules: four mathematically zero geometry blocks reveal roundoff.
The corrected graph passes all 891 real rays, the original 32 regression
models and another fixed 32 synthetic models from seed 13019, with
nonzero screen kernels kept in the synthetic group. The real maximum
whole-ray relative L2 error against the prior independent MP raw chain
is $9.15\times10^{-11}$; maximum block error is
$1.73\times10^{-10}$. All 64 synthetic comparisons pass, with maximum
whole-ray relative L2 $4.59\times10^{-14}$. The two retained real
insufficient-mass updates also pass value/directional checks. Ten
separate coefficient tests include the singular median, mass boundary,
strict amplitude cap and separated thin gaps; invalid input is explicitly
rejected. The known vanishing-gradient example remains unchanged.

The real batch takes three numerical branches: 877 interior solves,
11 tail formulas and three zero losses. No ray is dropped. On the
RTX 5090 Laptop GPU, the corrected 891-ray geometry/loss/raw-backward
batch has warm observed times $79.164$, $79.067$, $79.066$ ms;
the first cold observation is $1045.543$ ms and process peak allocated
memory is 153,932,288 bytes. These event timings exclude preparation,
host transfer and compilation. The initial version's separate stage
probe gives warm medians $5.54$ ms for fields, $72.70$ ms for loss/seeds
and $10.98$ ms for raw backward, identifying the serial-per-ray kernel
as a current cost. It is not a full-training acceleration claim.
Independent source/binary/archive and all-array audits pass. Earlier
confirmation serialization failures and the original 28/32 numerical
failure remain recoverable. Training, new real scenes, held-out images,
ground-truth improvements and novelty are still unverified.

\subsection{Warp-parallel evaluation of the same loss}
A separate CUDA version distributes each ray over a 32-lane warp, with
strided contributors, compensated local sums, double-double prefix
products and a fixed warp reduction tree. It retains the same loss,
independent root solve, caps and coefficient seeds, and uses the same
corrected raw graph. It has no silent serial fallback in its numerical
or timing results. The serial source and binary remain unchanged.

Against the previously saved independent MP raw references, both versions
pass all 891 real development rays, all 64 synthetic raw fixtures, the
two retained insufficient-mass updates, and the ten special coefficient
cases. Two additional rays have 513 and 1025 contributors. In each,
the prefix without the last primitive has final transmittance above
$1/2$, while the complete list has final transmittance below $1/2$;
an independent exact-rational check confirms these inequalities.
Both long rays pass the full-list MP coefficient comparison and report
the complete node count. Invalid zero-width input is explicitly rejected.
The maximum real whole-ray and parameter-block relative errors are
$9.15\times10^{-11}$ and $1.73\times10^{-10}$, respectively.

In one process on the same device and graph, the two versions alternate
execution order for ten paired rounds. The first two rounds per version
are retained as warmups; all eight later observations contribute to the
reported medians. Complete geometry/loss/raw-backward GPU event medians
are $79.1189$ ms for serial and $11.9931$ ms for warp, a ratio of
$0.151583$ (component speedup $6.60$). Loss/seed-only medians are
$72.7598$ and $5.9131$ ms. Allocation is included; setup, transfer and
compilation are separate. This is acceleration of this loss implementation
relative to its serial reference, not of the original renderer or an
entire training step. The observations are not confidence intervals or
independent training repetitions. The original correctness gates and
prespecified full-batch ratio threshold $0.95$ are both met.

An independent audit checks 111 sources, 17 saved artifacts, the actual
binary, all loss/gradient arrays, the frozen MP reference mapping, exact
long-ray mass witnesses and all alternating timing observations. An
earlier runner compiled the same warp kernel and finished timings but
failed while constructing a summary dictionary; its sources and failed
terminal state are retained. The corrected runner reuses that binary
and repeats the same fixed measurement protocol. Its field named
\texttt{compile\_wall\_seconds} records binary loading time in this
recovery run, and is not reported as compilation speed.
The warp version is adopted for the next research adapter. Live sampling,
full renderer visibility, zero-weight controls, optimizer updates,
held-out images and ground-truth quality remain to be validated.

\subsection{Live shared-parameter integration smoke}
The next diagnostic loads the actual 30k PLY models for scans 37, 55,
63, and uses camera 0000 in each. The original public renderer and the
private query/CSR captures agree bit for bit on RGB, alpha, normals,
depth, slope and radii. Actual query coordinates, amplitudes, accepted
IDs, source coefficients and raw values match the saved probe inputs.
Only sampled rows are packed to Float64 on GPU; their gradients scatter
back to the original shared Float32 GaussianModel parameters. The
corresponding independent MP loss gradients, summed by actual primitive
ID, agree within whole-array relative L2 errors $4.69\times10^{-8}$,
$2.57\times10^{-8}$ and $3.85\times10^{-8}$; all raw blocks pass.
No unintended appearance/color/SG regularizer gradient is observed.

Input compatibility requires an explicit qualification. The original
PLYs already contain 3431, 1908 and 2443 positive-infinity opacity logits,
respectively, with no opacity NaN or negative infinity. The initial
blanket finite-parameter check fails before any optimizer update. The
amended smoke preserves these exact entries as fixed opaque coordinates,
with sigmoid value one and exactly zero gradient; their masks and values
must remain unchanged. It does not clamp or repair the original files.
All remaining parameters, every gradient and all Adam states must be
finite. This is not a differentiability claim at an infinite logit, nor
evidence about the historical cause of those values or prior regressions.
An unrelated first scatter comparison also used a median-depth reference
instead of the history-loss reference; that oracle-selection failure is
retained and corrected by verifying the exact loss-reference mapping.

The sampled smoke batches contain 164, 162 and 167 rays, respectively,
all nonempty in these observations. Empty-list behavior is explicit:
the extended objective at $\Omega=0$ is $(\log4)^2=1.9218120557$,
with zero derivative for the fixed empty contributor graph. A real
empty ray is tested in scan37. A separate replay fixture combining its
36 saved nonempty probe rays and one empty ray preserves the denominator:
the mixed loss matches $(36L+(\log4)^2)/37$ within
$5.56\times10^{-17}$, its shared gradient has the expected $36/37$
factor within relative error $8.29\times10^{-8}$, and the all-empty
gradient is exactly zero. This fixture is not a new native-capture scene.

With original parameters restored between arms, two disabled repeats,
a computed zero-weight arm and a weight $0.01$ arm each take three
fresh-Adam RGB-only steps: 36 updates in total, without topology or
filter changes. The objective uses L1 and SSIM; the full NCC/normal
training recipe and resumed optimizer are not exercised. Zero-control
gradient differences pass the prespecified repeat-noise engineering
bound, not a bitwise or statistical equivalence claim. For example,
scan37's largest zero-control group difference is $1.0542\%$, versus
$0.6782\%$ in its disabled repeat. Maximum three-step zero-control
parameter relative differences are $2.79\times10^{-6}$,
$8.43\times10^{-6}$ and $4.27\times10^{-6}$. Optimizer bindings,
finite-state conditions and RNG preservation pass, and original
parameters and PLY hashes are restored/verified.

Across the positive arm's three pre-update observations, regularizer
losses change from $0.163931$ to $0.138692$, $0.244935$ to $0.224625$,
and $0.123730$ to $0.104279$. At the same final observations, the RGB
mixed losses are slightly higher than the disabled controls. This
supports an operative regularizer with a local tradeoff, not improved
image or reconstruction quality. Warm disabled-repeat synchronized
update-wall medians are $70.71$, $47.25$, $31.30$ ms; positive-arm
medians are $111.31$, $85.66$, $64.38$ ms. Positive capture alone takes
approximately $42.32$, $30.51$, $24.20$ ms and packing/preparation/loss
another $6.36$, $5.37$, $5.92$ ms. These smoke costs include finite and
binding checks and synchronization, but exclude artifact transfers/writes
and initial-state restoration. The isolated 12 ms loss result is not
the full adapter overhead. An independent audit checks all saved replay,
scatter, sample normalization, opaque masks, controls, inputs and step
statistics. Full training, held-out scenes/images and GT remain open.

\subsection{Original full-recipe integration diagnostic}
A separate scan55 diagnostic restores the strictly 42-view-trained
iteration6999 checkpoint and its runtime sidecar, including optimizer,
filter, densification statistics, camera queue and RNG. The canonical
complete-state digest agrees in four separate processes. The initial
676277 primitives include 752 existing positive-infinite opacity logits,
handled as the fixed opaque coordinates described above. The original
NCC and geometric losses have weights $0.6$ and $0.02$; the declared
depth-normal weight is zero. Actual train/test camera guards enforce
the 42/7 split for rendering, NCC neighbors and filter updates.

Two disabled repeats, computed zero weight, and weight $0.01$ each
complete six Adam calls, iterations7000--7005: 24 updates in total.
Original densification and filtering remain active; all arms have
677291 primitives after the first densification. Full gradients saved
immediately after the first backward precede that topology change.
For position, opacity, scale and rotation, zero-versus-disabled relative
gradient differences are respectively $1.58\times10^{-7}$,
$1.51\times10^{-7}$, $4.63\times10^{-5}$ and $3.01\times10^{-4}$.
Disabled-repeat differences are $1.79\times10^{-7}$,
$1.36\times10^{-7}$, $6.98\times10^{-5}$ and $2.48\times10^{-4}$.
All parameter groups pass the prespecified bound
$\max(10^{-5},5\,\text{repeat difference})$; this is an engineering
control, not statistical equivalence. Positive-weight first-gradient
changes in those four geometry groups are approximately $32.89\%$,
$15.25\%$, $14.83\%$ and $22.72\%$. Thus the regularizer contributes
to the real full-loss backward. This does not imply every contribution
survives original parameter replacement during densification.

Each weighted step samples128 uniform image pixels with a separate
generator; camera/sample sequences agree between weighted controls and
training RNG is unchanged. Finite-state and optimizer-binding checks
pass throughout, with the current opacity mask checked around each
ordinary Adam call. Independent artifact and gradient recomputation
passes. These camera-varying six-step losses do not establish a quality
trend. A fixed full-budget paired pilot subsequently tests official DTU
geometry, seven excluded images, capacity and complete resource cost;
its quality outcome is pending. No new quality or acceleration claim
is made from this diagnostic.

\subsection{Fixed-budget full pilot: a negative result}
Both scan55 branches subsequently restart from the identical common6999
state and run to30000, each with23000 further Adam calls. Final30000
includes backward, reporting and saving but no Adam. The control disables
the new term, and the candidate uses weight0.01. All16 orchestration
stages finish normally, including official-style mesh evaluation at
three sampling seeds and all42 TRAIN/seven TEST images. An independent
audit recomputes the serialized step records, actual final checkpoint
and runtime-sidecar digests, per-view means and fixed gates.

\begin{center}
\begin{tabular}{lrr}
\toprule
Metric & Control & Weight0.01\\
\midrule
Mean DTU overall (mm; lower better) & 0.3307163221 & 0.3315905070\\
TEST global RGB PSNR (dB) & 28.07059016 & 28.07511820\\
TEST SSIM & 0.899857896 & 0.900708420\\
TRAIN global RGB PSNR (dB) & 35.54241317 & 35.46370216\\
Gaussian count & 928975 & 931140\\
Full training process wall (s) & 2387.70 & 3151.18\\
\bottomrule
\end{tabular}
\end{center}
Overall worsens by0.26433\%; both directed-distance ratios also exceed1
(1.00209755 and1.00294454). Thus the prespecified geometry gates fail,
including the required overall ratio at most0.99. TEST PSNR changes by
only0.004528dB. Image noninferiority, capacity and resource guards pass,
but all gates were required. The observed full-process cost increases
by31.98\%, including full CSR and checks; this is resource accounting,
not a controlled speed benchmark. One exploratory shared-warmup pair
and correlated views do not establish significance or general harm.
The three evaluator seeds are not independent training repetitions.

A declared post-hoc, read-only matched-query diagnostic samples512
pixels on each of the42 training cameras for both final models. Among
18344 common nonempty rays, mean implemented loss decreases from
0.1776675410 to0.0931691978 (47.5598\%). Full-sample means decrease
from0.4284622530 to0.3498538008; their empty-ray constant contributions
are0.2746339494 and0.2689142706. All84 saved camera arrays and sample
identities are independently recomputed. This supports a loss/task
mismatch in this fitted pair, not a causal explanation or a replacement
quality criterion. Automatic weight search and expansion of this
regularizer are stopped; the mathematical and numerical results remain.

\paragraph{Relation to established objectives.}
Depth concentration already appears in2DGS, with its implemented squared
NDC-depth distortion detailed in AppendixA~\cite{twodgs}. Triangle
Splatting also uses distortion/normal terms and a compact triangular
window~\cite{trianglesplat}; the present Gaussian coefficient theorem
does not automatically cover that representation. Gaussian Wrapping
changes the attenuation/normal model and densification to address field
consistency~\cite{wrapping}. This targeted comparison is not an exhaustive
novelty search and establishes no superiority of the current loss.

\section{Certified quantile brackets}
Define the strictly increasing normalized component functions
$F_i(t)=-\log T_i(t)/a_i$. For $k>0$, set $r=\tfrac12 e^{-k^2/2}$ and
\[
 C_+(t)=\sum_i a_i\ind[p_i-k\sigma_i\leq t],
 \qquad
 C_-(t)=\sum_i a_i\ind[p_i+k\sigma_i\leq t].
\]
For weights $a_i>0$, let
$Q_b(z)=\inf\{t:\sum_i a_i\ind[z_i\leq t]\geq b\}$.
\begin{theorem}[Weighted-quantile certificate]
Suppose $A>h$ and choose
$0<r<\min\{h/A,1-h/A\}$. With $k=\sqrt{-2\log(2r)}$, the finite values
\[
 L=Q_{(h-rA)/(1-r)}(p-k\sigma),\qquad
 U=Q_{h/(1-r)}(p+k\sigma)
\]
satisfy $L\leq t^*\leq U$.
\end{theorem}
\begin{proof}
Convexity and the zero endpoint of $-\log(1-\alpha z)$ give
$-\log(1-\alpha z)\leq z[-\log(1-\alpha)]$ for $0\leq z\leq1$.
Thus $F_i(p_i-k\sigma_i)\leq r$ and, by branch symmetry,
$F_i(p_i+k\sigma_i)\geq1-r$. Splitting contributors according to the two
thresholds gives
\[
 (1-r)C_-(t)\leq H(t)\leq rA+(1-r)C_+(t).
\]
For $t<L$, the upper bound is strictly below $h$; continuity implies
$H(L)\leq h$. At $U$, the lower bound is at least $h$.
Strict monotonicity proves the result. The condition on $r$ places both
quantile levels strictly between $0$ and $A$, so both endpoints are finite.
\end{proof}
Once a numerical method preserves $H(\ell)\leq h\leq H(u)$, its midpoint
has absolute root error at most $(u-\ell)/2$. Floating-point endpoint checks
are still necessary. The current CPU choice
$r=\tfrac14\min\{h/A,1-h/A\}$ is conservative, with no optimality proof.
GGGS already identifies adaptive bracketing as a future direction; the present
certificate alone does not establish a new rendering algorithm.

\section{Differentiating the interface slope}
Let $S(\theta)=-T_t(t^*(\theta),\theta)=qH_t(t^*(\theta),\theta)$.
Where $H_t>0$ and the indicated second derivatives exist, the chain rule gives
\[
 \frac{dS}{d\theta}
 =q\left(H_{t\theta}-\frac{H_{tt}H_\theta}{H_t}\right).
\]
The root-motion term cannot in general be dropped. Moreover, $H_t>0$ alone
does not ensure that these second derivatives exist.

\begin{proposition}[A regular root with a slope cusp]
There are admissible two-primitive configurations for which $t^*(\theta)$ is
$C^1$ but $S(\theta)$ is not differentiable.
\end{proposition}
\begin{proof}
Take $q=1/2$, $p=(0,0.4)$, $\sigma=(0.1,0.4)$, $\alpha_2=0.8$,
$b=0.8e^{-1/2}$, and $\alpha_1=1-0.25/(1-b)$.
Then $T(0)=1/2$ and $H_t(0)>0$ due to the second component.
Vary only $p_1=\epsilon$. At $\epsilon=0$, $H_{p_1}=0$, so
$dt^*/d\epsilon=0$ and $t^*(\epsilon)=o(\epsilon)$.
The first component of $qH_t$ contributes
\[
 \frac{q\alpha_1}{2\sigma_1^2(1-\alpha_1)}
 |\epsilon|+o(|\epsilon|).
\]
The second component is smooth near $t=0$ and changes by $o(\epsilon)$.
The two one-sided derivatives of $S$ are therefore opposite nonzero values.
\end{proof}
Away from all contributor peaks, the displayed full derivative is valid.
At the cusp a chosen subgradient is an additional optimization convention,
not a classical derivative. Translating all one-dimensional peaks by the
same constant always translates the root and preserves its slope; this
provides a useful directional invariance check even where individual
partial derivatives do not exist.

\section{Local adjoint refinement and its domain of validity}
We next consider an approximate root $t_0$ and the continuous implicit
Jacobian $J_*=-H_\theta(t^*)/H_t(t^*)$. This is a Jacobian of the
underlying continuous root, not the exact derivative of a finite
subdivision program. Classical inverse quadratic interpolation and
Halley's rule are established numerical tools~\cite{brentdocs,dlmf}.
Approximation analyses of one-step differentiation also precede this
note~\cite{onestep}. The following bounds specialize standard local
arguments to the quantities needed by this implementation.

\begin{proposition}[Newton with a linearized implicit denominator]
Let $I$ be an interval containing $t_0,t_1,t^*$ and no contributor peak,
where $t_1=t_0+d_N$ and
\[
 d_N=\frac{T(t_0)-q}{T(t_0)H_t(t_0)},\qquad
 \widehat D=H_t(t_0)+H_{tt}(t_0)d_N.
\]
Assume on $I$ that $|T'|\geq m_F>0$, $|T''|\leq M_F$,
$H_t\geq m_H>0$, $|H_{tt}|\leq M_H$, $|H_{ttt}|\leq L_H$,
$\|H_\theta\|\leq B$, and $\|\partial_tH_\theta\|\leq N$.
Suppose also that $\widehat D\geq m_H/2$. Put $e=|t_0-t^*|$
and $e_1=|t_1-t^*|$. Then
\[
 e_1\leq\frac{M_F}{2m_F}e^2,
\]
and the refined estimate $\widehat J=-H_\theta(t_1)/\widehat D$ satisfies
\[
 \|\widehat J-J_*\|\leq
 \frac{2Ne_1}{m_H}+\frac{2B}{m_H^2}
 \left\{\frac{L_H(e+e_1)^2}{2}+M_He_1\right\}.
\]
\end{proposition}
\begin{proof}
Taylor expansion of $F=T-q$ between $t_0$ and $t^*$, divided by
$F'(t_0)$, proves the Newton bound. The denominator remainder obeys
$|\widehat D-H_t(t_1)|\leq L_Hd_N^2/2$ and
$|H_t(t_1)-H_t(t^*)|\leq M_He_1$. The numerator difference is at
most $Ne_1$. Subtract the two ratios, use both denominator lower bounds,
and note that $|d_N|\leq e+e_1$.
\end{proof}

For $q=1/2$, write $z_i=(t-p_i)/\sigma_i$ and $R=qH_t$.
On either smooth branch, direct differentiation yields
\begin{align}
 R'&=\frac14\sum_i\frac{g_i\sgn(z_i)}{\sigma_i^2}
       \frac{1-g_i-z_i^2}{(1-g_i)^2},\\
 R''&=\frac14\sum_i\frac{g_i|z_i|}{\sigma_i^3}
       \frac{(1+g_i)z_i^2-3(1-g_i)}{(1-g_i)^3}.
\end{align}
With $\rho=(T(t_0)-q)/T(t_0)$, the classical Halley step is
\[
 d_H=\frac{\rho}{H_t(t_0)\{1-\tfrac12\rho
             [1-H_{tt}(t_0)/H_t(t_0)^2]\}},\qquad
 \widehat R=R+R'd_H+\tfrac12R''d_H^2.
\]

\begin{proposition}[Conditional cubic Jacobian accuracy]
Fix an admissible parameter configuration and a simple root $t^*$
away from every contributor peak. Suppose $T$ is $C^3$, $H$ is $C^4$,
$H_\theta$ is Lipschitz in $t$, and these derivatives are bounded on
a neighborhood with $|T'|,H_t$ bounded below by positive constants.
For $t_0$ sufficiently close to $t^*$, let $t_H=t_0+d_H$ and define
$\widehat J_H=-qH_\theta(t_H)/\widehat R$. Then
\[
 |t_H-t^*|=O(|t_0-t^*|^3),\qquad
 \|\widehat J_H-J_*\|=O(|t_0-t^*|^3).
\]
The constants and neighborhood depend on the fixed configuration.
\end{proposition}
\begin{proof}
Write $f_1=F'(t^*)\ne0$, $a=F''(t^*)/(2f_1)$ and
$b=F'''(t^*)/(6f_1)$. For signed $u=t_0-t^*$,
$F(t_0)/f_1=u+au^2+bu^3+o(u^3)$,
$F'(t_0)/f_1=1+2au+3bu^2+o(u^2)$ and
$F''(t_0)/f_1=2a+6bu+o(u)$. Substitution into
$t_H=t_0-2FF'/(2(F')^2-FF'')$ gives
$t_H-t^*=(a^2-b)u^3+o(u^3)$. The denominator is nonzero
for $u$ sufficiently small. Since $R=qH_t$ is $C^3$,
$\widehat R-R(t_H)=O(|d_H|^3)=O(|u|^3)$.
The true denominator at $t_H$ differs from that at $t^*$ by
$O(|u|^3)$, as does the numerator by its Lipschitz condition.
A ratio difference with positive denominator lower bounds proves the claim.
\end{proof}

\begin{remark}[Loss gradients and finite precision]
If the forward loss is still evaluated at $t_0$, its upstream factor
is $\ell'(t_0)$. Even with a quadratically or cubically accurate
implicit Jacobian, $\ell'(t_0)-\ell'(t^*)$ is generally $O(e)$.
The complete loss gradient therefore need not have the same higher order.
For example, a single regular primitive with $\alpha=0.9$, $p=1$,
$\sigma=0.1$, $q=1/2$ has
$t^*=p-\sigma\sqrt{-2\log[(1-q^2)/\alpha]}$ and $dt^*/dp=1$.
For $\ell(t)=t^2/2$, a high-order Jacobian multiplied by the coarse
upstream $t_0=t^*+e$ still has leading loss-gradient error $e$.
\texttt{verify\_adjoint\_loss\_limit.py} independently confirms this order.
Neither proposition supplies uniform constants as $H_t\to0$,
$\alpha_i\to1$, $\sigma_i\to0$, or a peak is approached.
A CUDA step guard such as $|d_H|<\tfrac12\min_i|t_0-p_i|$
keeps its accepted segment away from peaks, but does not certify the
true root, the lower denominator bounds, or floating-point roundoff.
\end{remark}

The engineering prototype accumulates $T,R',R''$ in the existing
backward denominator traversal and evaluates numerators at the refined
local root in its existing second traversal. Thus it introduces arithmetic
and registers but no additional traversal. Whether this saves time or
preserves a training trajectory requires measurement.

\section{A local normal-normalization derivative discrepancy}
The Gaussian preprocessing forward normalizes an unnormalized camera
normal $c$ to $n=c/\|c\|$. The inspected local and official upstream
backward sources~\cite{gggscode} instead compute the reciprocal length
from the already normalized vector. For nonzero $c$ and upstream $g$,
the correct local VJP is
\[
 \nabla_c L=\frac{g-n(n^\top g)}{\|c\|}.
\]
This follows directly from $dn=(I-nn^\top)dc/\|c\|$.
Replacing the denominator by $\|n\|=1$ generally changes the derivative.
For a nondegenerate example take camera coordinates $(0,0,1)$ and
$P=\left(\begin{smallmatrix}2&0&x\\0&1&0\\x&0&1\end{smallmatrix}\right)$,
which is positive definite near $x=1$. The affine preprocessing gives
$c=(-x,0,-1)$. For $L=n_x$ the true derivative is
$-(1+x^2)^{-3/2}$, whereas the inspected local expression gives
$-(1+x^2)^{-1}$; at $x=1$ these are $-1/(2\sqrt2)$ and $-1/2$.
Symbolic, float64 autograd and directional finite-difference checks agree.
This is a local discrepancy, not a complete world-gradient audit or
a reconstruction improvement. A geometry-only loss has zero upstream
normal gradient, so this factor cannot explain its numerical variability.

On the positive-definite affine branch this normal has a simpler form.
Let $d$ be the unit center-view direction and $P$ the camera precision.
Writing $r=(u,v,1)$, $\rho^2=1+u^2+v^2$, $h=Pr$, $b=r^\top h$,
the two plane components in the implementation are
$\eta_i=\rho^2h_i/b-r_i$ for $i=1,2$.
Substitution into its camera transformation, with factor $z/\rho$,
gives $c=-\rho h/b$. Since $d=r/\rho$, it follows that
\[
 c=-\frac{Pd}{d^\top Pd},\qquad
 n=-\frac{Pd}{\|Pd\|},\qquad
 \|c\|=\frac{1}{|n^\top d|}\ge1.
\]
The last inequality follows from Cauchy--Schwarz and $\|d\|=1$;
positive definiteness makes $d^\top Pd>0$.
Thus the omitted inverse length is exactly the absolute incidence
cosine within this affine primitive model. Its direct differential is
\[
 dn=-\frac{(I-nn^\top)(dP\,d+P\,dd)}{\|Pd\|}.
\]
Any active FoV clipping must also be differentiated.
Uniform positive scaling of $P$ leaves $n$ invariant; consequently
the sum of the three log-scale derivatives of a normal-only loss
vanishes. This invariant does not apply to the complete RGB or
geometry loss. Three symbolic component identities and 64 independent
float64 world-VJP computational-graph comparisons pass. These algebra
checks alone supply no CUDA derivative, runtime or reconstruction-quality
guarantee; the subsequent isolated CUDA audit is reported below.

\section{Two further preprocessing chain rules}
These statements concern the positive-scale branch of the actual
affine primitive preprocessing, separately from implicit differentiation
of the median root. They establish derivatives relative to the implemented
forward model, not a new surface representation.

\paragraph{Ray precision and its center derivative.}
Let $P$ be the symmetric positive-definite camera precision matrix,
$r=(u,v,1)^\top$, $L=r^\top r$, $h=Pr$, $B=r^\top Pr$, and let
$\kappa=\sqrt{B/L}$, named \texttt{rsigmat} in the shader.
On this smooth positive branch,
\[
 \frac{\partial\kappa}{\partial u}
 =\frac{h_1}{\kappa L}-\frac{\kappa u}{L},\qquad
 \frac{\partial\kappa}{\partial v}
 =\frac{h_2}{\kappa L}-\frac{\kappa v}{L}.
\]
Indeed $\partial_u B=2h_1$ and $\partial_u L=2u$;
the quotient and square-root rules give the first identity, and the
second follows identically. The existing local backward includes
the second term of each identity but omits the numerator term. Its
precision-matrix derivative elsewhere does not provide the missing
center/viewport derivative. For $P=I$, $\kappa=1$ is constant, whereas
the existing rule gives $-2/5$ at $(u,v)=(1/2,0)$ in the $u$ direction.
This is a non-singular positive-definite counterexample. The clamp
and world-to-camera transformations must still be chained for world
centers; the proposed correction retains those existing transformations.

\paragraph{Opacity filtering coefficient.}
Let $D_0,D_1>0$ be the unclamped projected-covariance determinants
before and after adding the image-space filtering kernel, and
$\eta=\sqrt{D_0/D_1}$. Direct differentiation yields
\[
 \partial_{D_0}\eta=\frac{1}{2\eta D_1},\qquad
 \partial_{D_1}\eta=-\frac{\eta}{2D_1}
 =-(\partial_{D_0}\eta)\eta^2.
\]
The existing backward multiplies the first derivative by $-\eta$
instead of $-\eta^2$. The opt-in correction changes that factor;
it retains the existing coefficient epsilon and determinant floors,
so it is not a claim of exact differentiation at every stabilized
or clamped branch. With filtering kernel zero, $D_0=D_1$ and
$\eta=1$; this particular factor correction has no effect in ideal
arithmetic. All three real models audited here use kernel zero;
the nonzero-kernel discrepancy is demonstrated by separate synthetic
cases. The same source patterns were observed in official main on
2026-10-12~\cite{gggscode}; upstream history and author intent were not audited.

\paragraph{Artificial angular term and non-integrability.}
Fix $P$, set the upstream of $\kappa$ to one and all other output
upstreams to zero, and work inside the FoV clamp. The existing viewport
field is $g=(-\kappa u/L,-\kappa v/L)$. Differentiating it and using the
preceding exact identities gives
\[
 \partial_u g_v-\partial_v g_u
 =\frac{u h_2-v h_1}{\kappa L^2}.
\]
The derivatives of $1/L$ cancel; the remaining numerator is
$u\partial_v\kappa-v\partial_u\kappa$, divided by $L$.
For $P=\operatorname{diag}(1,2,3)$ and $u=v=1/2$, this curl is
$\sqrt{10}/45\ne0$. Consequently this component of the existing
backward cannot be the gradient of any $C^2$ local scalar potential
in a neighborhood of that point, by equality of mixed partials.
This statement concerns a fixed-upstream component, not a proposed
surrogate for the complete image or reconstruction loss.

For the separate isotropic case $P=\lambda I$, $\kappa=\sqrt\lambda$
is independent of the center direction, but the existing center field
in camera coordinates $t=(x,y,z)$, $z>0$, equals
\[
 g_t=\kappa\left(\frac{e_z}{z}-\frac{t}{\|t\|^2}\right)
 =\nabla_t[\kappa\log\cos\theta],\qquad
 \|g_t\|=\frac{\kappa\tan\theta}{\|t\|},
\]
where $\theta$ is the angle from the positive camera axis. Indeed
$\log\cos\theta=\log z-\log\|t\|$; direct differentiation gives
the vector and its squared norm. Thus a spurious angular term appears
in this isolated derivative although the true directional derivative
of $\kappa$ vanishes. Its effect in an actual loss depends on the
upstream sign and magnitude. No conclusion about the complete center
gradient vanishing, scene rotation invariance, or reconstruction bias
follows without additional analysis.
An independent script verifies the symbolic curl and angular identities,
64 float64 AD cases, and closed-loop quadrature: counterclockwise
circulation per unit area tends to $0.0702728369$ for the stated
anisotropic example, while the corrected field's observed absolute
circulation is below $3.5\times10^{-18}$. These are standard calculus
checks supporting the implementation diagnosis, not a novelty claim.

\paragraph{Independent isolated CUDA world VJPs.}
A thin probe calls the actual forward and backward preprocessing
functions with fixed normal, ray-plane and conic upstreams. A separate
float64 PyTorch graph uses camera covariance/precision matrices and
the simplified normal $-Pr/\|Pr\|$, with three directional finite-difference
step sizes checking that graph. Packed-conic off-diagonal upstreams
respect the shader's factor-of-two convention. Scale activation and
quaternion normalization are included in the pullback.
The error rule is $\|g_{\rm CUDA}-g_{\rm AD}\|_2/
\max(\|g_{\rm AD}\|_2,10^{-5})\le10^{-3}$ per parameter block.
Forty-eight synthetic primitives cover interior and active FoV-clamp
domains and a nonzero image-space kernel. A second audit selects 72
projected in-FoV primitives across frozen scan37/55/63, including
24 with filtered axis ratio above 10. Projection membership is not
visibility. The combined FP64 correction passes all 1344 synthetic
and 1440 real parameter-block checks, including structural zeros;
the largest errors are about $6.30\times10^{-6}$ and
$1.20\times10^{-4}$, respectively. The real first-order reference
also agrees with the smallest finite-difference step within
$1.18\times10^{-6}$ on the stated scaled rule.
The combined FP32 version fails the same rule. In particular, some
zero-reference opacity-coefficient checks use the absolute floor;
their scaled errors must not be described as relative errors of a
nonzero true gradient. These tests do not certify extreme anisotropy,
the eigendecomposition fallback, complete rasterization, hard masks,
optimization trajectories, quality or novelty. Full-backend function
bytes are checked against the audited probe before training.

\section{A direct inverse-covariance factor VJP}
This separate numerical-stability diagnostic does not change the median
root model. Let $S=\operatorname{diag}(s_1,s_2,s_3)$ with $s_i>0$,
let $R$ be orthogonal, and write
$\Sigma=R^\top S^2R$, $P=\Sigma^{-1}=N^\top N$, $N=S^{-1}R$.
For the local scalar loss $L=\langle G,P\rangle_F$ with fixed upstream
matrix $G$, the Frobenius differential gives
\[
 G_N=N(G+G^\top),\qquad
 \frac{\partial L}{\partial s_i}
 =-s_i^{-2}\langle R_{i,:},(G_N)_{i,:}\rangle,\qquad
 G_R=S^{-1}G_N.
\]
Indeed $dP=dN^\top N+N^\top dN$ and
$dN=-S^{-1}(dS)S^{-1}R+S^{-1}dR$; substituting these two identities
proves the displayed formula. On the orthogonal-rotation manifold it
agrees with the inverse-matrix chain $G_\Sigma=-PGP$ followed by
differentiation of $\Sigma$. For a normalized quaternion one must
also differentiate its normalization; an arbitrary nonorthogonal
matrix $R$ does not satisfy the inverse identity.
The factor form removes an unnecessary inverse-to-covariance round
trip but provides no uniform conditioning or floating-point guarantee.
Synthetic symbolic and float64 autograd checks pass; a real-model
GPU diagnostic still exhibits substantial full-loss world-parameter
gradient variability after this local change. It is therefore not
adopted as a verified stability or quality improvement. These are
standard matrix-calculus identities, not a novelty claim.
In a separate precision diagnostic, FP64 intermediates in the local
world-parameter backward preprocessing substantially reduce repeat
variation on two selected development cameras. An expanded 18-camera,
162-group complete-loss check then passes the original 0.1\% rule,
with maximum drift $3.3312\times10^{-5}$ between FP64 sampling
batches 1 and 2. A selective-precision variant with filtered axis
ratio above 10 also passes, with maximum drift $1.0293\times10^{-4}$.
However, both variants fail to improve the six-camera frozen complete
loss/backward latency relative to the original float32 control.
The FP64 reference shares the analytic formulas and float32 contributor
accumulation; repeatability is not an independent correctness oracle,
and none of these checks establishes reconstruction quality.

\section{Orthogonality defects and anisotropic covariance round trips}
Here $R$ denotes the conventional column-vector rotation matrix; it is
the transpose of the preceding section's row-factor convention. Define
$D=\operatorname{diag}(s_1,s_2,s_3)$, $s_i>0$, and
\[
 \Sigma=RD^2R^\top,\qquad P=RD^{-2}R^\top,\qquad
 Q=R^\top R-I,\qquad \delta=\|Q\|_2,\quad
 a=s_{\max}/s_{\min}.
\]
Neither definition requires $R$ to be exactly orthogonal. Direct
multiplication gives the exact identity
\[
 P\Sigma-I=(RR^\top-I)+RD^{-2}QD^2R^\top.
\]
For square $R$, $R^\top R$ and $RR^\top$ have the same eigenvalues,
so $\|RR^\top-I\|_2=\delta$ and $\|R\|_2^2\le1+\delta$.
Submultiplicativity therefore proves
\begin{equation}\label{eq:roundtrip-defect}
 \|P\Sigma-I\|_2\le
 \delta\bigl[1+(1+\delta)a^2\bigr].
\end{equation}
This sufficient upper bound identifies a quadratic dependence on axis
ratio. It is not a lower bound, a prediction that every anisotropic
primitive fails, or a bound on the full renderer's gradient error.

For the standard quaternion polynomial, let $q=(w,v)$,
$K=[v]_\times$, and $R(q)=I+2wK+2K^2$. Without assuming unit norm,
\[
 R^\top R-I=4(\|q\|^2-1)(\|v\|^2I-vv^\top),\qquad
 \delta=4|\|q\|^2-1|\,\|v\|^2.
\]
To verify this, use $K^\top=-K$, $K^2=vv^\top-\|v\|^2I$ and
$K^4=-\|v\|^2K^2$ in $(I-2wK+2K^2)(I+2wK+2K^2)$.
All nine polynomial components are also checked exactly with SymPy.
A quaternion normalized in float32 need not have exactly unit norm
when its stored components are subsequently interpreted in float64.
Increasing arithmetic precision after that node does not restore
the missing orthogonality identity. Applying a covariance inverse
round trip to differentiate the factor-defined $P$ can consequently
amplify this discrepancy. This statement concerns the actual stored
node values, rather than an error in exact unit-quaternion algebra.

\paragraph{The quadratic axis-ratio order cannot be removed.}
For fixed $\varepsilon>0$, take
$q=(0,\sqrt{(1+\varepsilon)/2},\sqrt{(1+\varepsilon)/2},0)$ and
$D=\operatorname{diag}(1,a,1)$, $a\ge1$. Its rotation polynomial is
\[
 R=\begin{pmatrix}
 -\varepsilon&1+\varepsilon&0\\
 1+\varepsilon&-\varepsilon&0\\
 0&0&-1-2\varepsilon
 \end{pmatrix},\qquad \delta=4\varepsilon(1+\varepsilon).
\]
Writing $r_i=R_{:,i}$, the exact defect identity implies
$(P\Sigma-I)/(\delta a^2)\to-r_1r_2^\top/2$.
The norm of a rank-one matrix is the product of the vector norms;
continuity therefore gives
\[
 \lim_{a\to\infty}\frac{\|P\Sigma-I\|_2}{\delta a^2}
 =\frac{\|r_1\|_2\|r_2\|_2}{2}
 =\tfrac12+\varepsilon+\varepsilon^2>0.
\]
Thus no finite axis-ratio-independent constant $C$ can bound this
residual by $C\delta$ over the stated family. This proves sharpness
of the quadratic order, not sharpness of the upper-bound constant
or a corresponding claim for loss gradients. The family is an exact
off-unit-quaternion counterexample; it is not asserted that each
member is produced by a float32 normalization kernel. A separate
symbolic limit check passes, with float64 corroboration at
$\varepsilon=2^{-24}$ and six axis ratios from 10 to $10^6$.

For any real $R$ and positive scales, differentiating the factor
definition itself avoids this orthogonality assumption. With
$G_s=(G+G^\top)/2$ and $\ell_i=\log s_i$, the differential of
$\langle G,RD^{-2}R^\top\rangle_F$ gives
\[
 G_R=2G_sRD^{-2},\qquad
 \frac{\partial L}{\partial\ell_i}
 =-2s_i^{-2}(R^\top G_sR)_{ii}.
\]
The external raw-quaternion normalization must still be pulled back.
These are standard matrix-calculus identities, with no novelty claim
or guarantee of uniform floating-point conditioning.

\paragraph{Stress failure, replacement and scope.}
A pre-quality extension selects 44 extreme-axis primitives across
the same three models, including both in-FoV and front-facing cases
outside FoV. The largest axis ratio is about 34819. The formula-fixed
FP64 covariance-round-trip branch fails 12 normal, 3 plane-slope and
28 ray-precision parameter blocks out of 176 per output. The worst
ray-precision error is 47.294 on the original scaled rule, with
nonzero reference norm 0.364575. Four extreme forward components
also fail their original gate; those failures remain recorded.
The observed largest $\|P\Sigma-I\|_2$ is 28.0014. The bound in
\eqref{eq:roundtrip-defect} is checked on all 164 saved primitives
within a stated $10^{-5}$ numerical matrix-product margin; the proof
above is in exact arithmetic. Sixty-four independent nonorthogonal
factor-graph AD examples agree within $1.43\times10^{-15}$.

Combining the direct factor VJP, the three formula corrections and
FP64 local intermediates passes all 880 stress checks and all 2784
regular real/synthetic checks without excluding inputs. The maximum
scaled errors are $1.68\times10^{-5}$ and $1.20\times10^{-4}$,
respectively. This combined branch does not change forward rendering,
and therefore does not repair the four retained forward failures.
It does not supersede the earlier negative float32 full-loss diagnostic.
The unstarted covariance-round-trip candidate was withdrawn; its
completed, preselected FP64 control is retained. A separate replacement
protocol preserves the original quality gates and requires a new
20-update full-backend smoke, which passes with the same initial state,
first camera and first forward losses. The replacement 30k pilot has
completed with the original forward unchanged. Its mean overall improves
by $0.001896019$ mm, or about $0.60186\%$, and mean s2d improves by
$0.002958633$ mm. The preregistered $1\%$ overall-improvement requirement
is not met; the geometric adoption verdict remains failure.
Mean TRAIN-view PSNR rises by $0.532053$ dB and SSIM by $0.00312850$,
while the Gaussian count rises from 943641 to 967292 ($2.50636\%$).
Thirty-five of 49 training views improve in PSNR; 14 worsen.
Capacity may mediate part of the change. One trajectory pair, correlated
training images and three mesh-evaluator seeds do not establish
significance, held-out quality, feature-level causation or novelty.
The three evaluator seeds repeat mesh sampling rather than training.

\paragraph{Related implementation analyses.}
A subsequent excluded-RGB pilot uses a fresh 42-train/7-test warmup,
not the earlier all-view checkpoint. Both branches complete all 19
training/evaluation stages. TEST PSNR changes from $28.032109$ to
$28.122184$ dB ($+0.090075$), below the predeclared $+0.30$ dB gate;
TEST SSIM changes by $-0.000304571$, also failing its gate.
Mean overall changes from $0.332269708$ to $0.331238519$ mm, while the
Gaussian count grows from $929179$ to $957753$ ($3.07519\%$).
The new RGB pilot therefore fails, independently of the preserved
failure of the original geometry pilot. TRAIN PSNR improves by
$0.421708$ dB; it must not be substituted for the excluded-view result.
Actual camera-pool guards exclude test RGB from training, NCC targets
and filter lists. Alignment now matches noncontiguous training image
IDs to their corresponding official calibration centers; a first-$N$
mapping would be invalid after holding out cameras. Existing SfM
initialization may use all views, and the scene was previously used
for research, so this is not an untouched confirmatory dataset.
There is no independent training replication or significance claim.

Faster-GS~\cite{fastergs} already studies numerical stability and
gradient approximations in Gaussian splatting. Its Appendix C.2
discusses non-analytical opacity-compensation gradients and offers
both a detached and an optional analytical derivative. Therefore
discovering a formula discrepancy alone establishes neither novelty
nor improved optimization. Our present quaternion/precision-factor
diagnosis needs its own source comparison and real-scene validation;
a targeted literature search is not a novelty certificate.

\section{A tight conditional log-scale VJP defect}
The round-trip residual alone is not a bound on a loss gradient.
We can, however, characterize a particular parameter block exactly.
Hold $R\in\R^{3\times3}$, a common symmetric precision upstream $H$,
and any filter fixed. Let $D=\operatorname{diag}(d_1,d_2,d_3)>0$,
$\ell_i=\log d_i$, $\Sigma=RD^2R^\top$, and $P=RD^{-2}R^\top$.
If $r_i$ is column $i$ of $R$, the actual factor graph has
\[
 g_i=-2d_i^{-2}r_i^\top H r_i.
\]
Using the inverse-chain surrogate $-PHP$ as a covariance upstream instead
gives $\widehat g_i=-2d_i^2(Pr_i)^\top H(Pr_i)$.
This comparison conditions on an identical $H$; it does not encompass
other chain-rule discrepancies or complete-renderer upstream changes.

\begin{proposition}[Exact defect and tight upstream norm]
Define $u_i=r_i/d_i$, $v_i=d_iPr_i$, $B_i=v_iv_i^\top-u_iu_i^\top$,
and $M_{ij}=\langle B_i,B_j\rangle_F$. Then
\[
 \widehat g_i-g_i=-2\langle H,B_i\rangle_F,
 \qquad
 \sup_{H=H^\top,\,\|H\|_F\le1}
 \|\widehat g-g\|_2=2\sqrt{\lambda_{\max}(M)}.
\]
For $Q=R^\top R-I$, the difference $w_i=v_i-u_i$ equals
$d_iRD^{-2}Qe_i$, and
$B_i=u_iw_i^\top+w_iu_i^\top+w_iw_i^\top$.
\end{proposition}
\begin{proof}
Differentiate $P$ directly with respect to $\ell_i$ to obtain
$-2d_i^{-2}r_ir_i^\top$. Differentiating $\Sigma$ gives
$2d_i^2r_ir_i^\top$, which paired with $-PHP$ yields the surrogate.
Their difference is the stated Frobenius inner product. The formula
for $w_i$ follows by substituting $R^\top R=I+Q$.
Let $K$ have rows $\operatorname{vec}(B_i)^\top$. Then $M=KK^\top$
and the error map is $-2K\operatorname{vec}(H)$.
If $\lambda=\lambda_{\max}(M)>0$ with unit eigenvector $z$, the symmetric
matrix $H_*=(\sum_i z_iB_i)/\sqrt\lambda$ has Frobenius norm one and
attains error norm $2\sqrt\lambda$. If $M=0$, the error vanishes.
\end{proof}
For fixed filtered scales $d_i^2=\exp(2\theta_i)+f_i^2$, the additional
raw-log-scale pullback is diagonal, $J_{ii}=\exp(2\theta_i)/d_i^2$.
The corresponding tight conditional norm is $2\sqrt{\lambda_{\max}(JMJ)}$.
This treats $f_i$ as fixed and excludes its potential separate dependence.

The non-unit quaternion family in the preceding section also illustrates
why an inverse residual and a relative VJP error are distinct. Set
\[
 R=\begin{pmatrix}-\varepsilon&1+\varepsilon&0\\
 1+\varepsilon&-\varepsilon&0\\0&0&-1-2\varepsilon\end{pmatrix},
 \quad D=\operatorname{diag}(1,a,1),\quad H=I,
\]
with fixed $\varepsilon>0$, $a\ge1$,
$\delta=4\varepsilon(1+\varepsilon)$ and $m=1+2\varepsilon+2\varepsilon^2$.
Direct expansion gives
\[
 \frac{\widehat g_2}{g_2}
 =m^2+\frac{\delta^2a^2}{2}+\frac{\delta^2a^4}{4},
 \qquad
 \lim_{a\to\infty}
 \frac{\widehat g_2/g_2-1}{\delta^2a^4}=\frac14.
\]
Here $g_2=-2m/a^2\ne0$. Replacing $H$ by $I/\sqrt3$ normalizes the
upstream without changing this ratio. Thus a relative component error
can grow quartically in this fixed-defect family. This is neither a
claim that all family members arise from float32 normalization nor a
lower bound on any complete training gradient.

Exact symbolic identities and the limit are checked in
\texttt{precision\_scale\_vjp\_defect\_v4.json}. Sixty-four independent
float64 automatic-differentiation cases verify both factor and surrogate
graphs, with maximum scaled discrepancy $9.285\times10^{-15}$;
constructed unit-norm upstreams attain the bound to $5.492\times10^{-16}$.
All 164 saved inputs were also retained. Under the particular trace
upstream $H=I$, their maximum whole-vector relative error is only
$7.045\times10^{-7}$, although one component ratio reaches $785.0793$.
The aggregate and component quantities therefore must not be conflated.
The trace upstream is not asserted to match the real training upstream.
Initial symbolic-development failures are archived separately.

\section{Projected Gram determinants and forward cancellation}
The four retained extreme forward failures all concern the conic block.
Let $a,b\in\R^3$ be the two projected columns of the scaled rotation
factor, and let $C_0=[a\ b]^\top[a\ b]$. Lagrange's identity gives
\[
 \det C_0=\|a\times b\|_2^2,
 \qquad
 \det(C_0+kI)=\|a\times b\|_2^2+
 k(\|a\|_2^2+\|b\|_2^2)+k^2,\quad k\ge0.
\]
The first equality follows by expanding
$(a^\top a)(b^\top b)-(a^\top b)^2$; the second follows by
expanding the two-by-two determinant after adding $kI$.
For $c=a\times b$, its untruncated derivatives are
$\nabla_a\det(C_0+kI)=2b\times c+2ka$ and
$\nabla_b\det(C_0+kI)=2c\times a+2kb$.
Thus the exact differentiable map is unchanged away from existing
floor and mask boundaries. The factor expression avoids the final
subtraction of two nearly equal covariance products, but can still
have cancellation within cross products, input-rounding sensitivity,
underflow or overflow. No uniform stability theorem is claimed.

A separate float32 forward-only probe replays all 44 stress and 120
regular real/synthetic inputs. On the original $10^{-4}$ forward
scaled rule with denominator floor $10^{-5}$, the conic failures
fall from four to zero and the largest error falls from
$2.90671\times10^{-3}$ to $2.76635\times10^{-5}$.
All six recorded output blocks pass on all 164 cases, and the ray
and normal outputs remain bitwise identical. A symbolic expansion
checks the Gram identity exactly. This experiment is independent of
the running factor64 quality candidate, whose original forward is
retained. The null-scale covariance fallback, complete-renderer
derivative correctness, training quality and latency are outside
this forward-only check. These are standard Cauchy--Binet/Lagrange
identities used as a numerical implementation diagnostic, not new
algebra.

\section{Numerical checks and limitations}
Independent scripts accompany the derivations. They are error-detection
checks, not substitutes for the proofs or real reconstruction experiments.
\begin{itemize}
\item \texttt{verify\_median\_solver.py}: 296 finite-root mixtures pass the
bracket checks against an independent Brent solver; maximum root error is
approximately $9.93\times10^{-10}$.
\item \texttt{verify\_root\_regularity.py}: a product-form transmittance and
Brent solver confirm the coincident-mixture expansion for 12 perturbations.
At step $10^{-7}$ the slope-cusp one-sided differences are
$-26.47748$ and $+26.47749$, versus the analytic magnitude $26.47755$.
\item \texttt{audit\_field\_slope.py}: moving-root implicit derivatives agree
with re-solved finite differences within $5.04\times10^{-6}$ in four
regular synthetic cases. The existing, explicitly documented field-occupancy
STE fails the one-dimensional common-peak translation check. This does not
establish the cause of past quality regressions.
\end{itemize}
The original hybrid Newton prototype uses a median of 33 log-transmittance evaluations
versus 32 for bisection, excluding sorting; this supplies no acceleration
evidence. A revised variant verifies a small local bracket near Newton convergence,
rather than continuing to shrink a distant opposite endpoint. In the same 296
random mixtures, it uses a median of 12 ordinary field evaluations, with 91
additional 80-digit checks across the family for ambiguous residual signs.
All 305 test cases pass, including additional near-threshold cases checked
against high-precision single-primitive formulas. The first float64-only
variant failed a flat-root endpoint check; that failure is retained.
These counts exclude sorting and unequal precision costs and are not a speed claim.
Sampled CPU/CUDA forward checks now pass on three frozen models with
1024 rays each, within the declared depth/alpha tolerances and with no
validity disagreements. Simple four-round subdivision lowers frozen
forward latency by about 14\% and preserves extracted meshes within
the preset engineering tolerance, but fails the 0.1\% parameter-gradient
drift rule. Inverse quadratic interpolation improves typical depths but
increases the maximum observed gradient drift to about 4.63\%.
The first local Newton adjoint refinement passes 21 of 24 parameter
groups; the Halley/quadratic-denominator variant passes all 24, with
maximum relative L2 drift $8.062\times10^{-4}$.
On six frozen development cameras its full-image RGB-plus-depth forward
and backward latency falls by about 6.5--8.7\%, excluding optimizer and
NCC work. These measurements do not establish full-training performance,
independent world-parameter derivative correctness, or novelty.
The scan55 shared-state 30k training pilot has now completed. Across
fixed evaluator seeds 0, 1, 2, the mean overall difference is
$+0.000424087$ mm and the mean s2d difference is $+0.000908171$ mm;
both are worse. Training-view PSNR increases by $0.066834$ dB.
The preset geometric adoption criteria are not met. Subsequent CUDA
training remains nondeterministic despite identical initial state;
one same-state control repeat has also completed. Its mean overall
difference from the original control is $+0.001001144$ mm, larger
than the original candidate contrast. A single paired trajectory
therefore cannot isolate the method effect. This repeat is not an
independent-seed or significance experiment, and does not replace
the original control or reverse the failed adoption verdict.

An independent product-form CPU float64 autograd oracle agrees with
the analytic effective-ray Jacobian within $4.12\times10^{-16}$
relative L2 on 123 previously used development rays. Re-solved
directional finite differences on 12 fixed rays also pass.
The local high-order guard accepts 122 rays, whose largest idealized
refined Jacobian error is $3.18\times10^{-5}$. It rejects one
peak-crossing ray; its coarse-root fallback retains relative error
$4.66\times10^{-3}$. Aggregated world-gradient agreement therefore
does not establish a bound for every ray. Float32 projection differences
also limit direct comparison to this CPU oracle. Independent multi-scene
training and held-out image evaluation remain necessary.
A peak-containing-cell variant restores subdivision five selectively.
It passes the 24 frozen gradient-group checks, but its six-camera
forward/backward latency is about 0.42--1.15\% worse than control;
the current implementation is rejected as an efficiency improvement.

The theoretical overlap with oriented Gaussian occupancy and surface
extraction~\cite{wrapping} also requires a broader novelty review.
Neither implicit differentiation nor adaptive bracketing is claimed as new.
The research workflow used Scientific Agent Skills~\cite{skills}; this
software citation is not evidence for the mathematical claims.

\begin{thebibliography}{99}
\bibitem{htgs}
F. Hahlbohm et al. \emph{Efficient Perspective-Correct 3D Gaussian
Splatting Using Hybrid Transparency}. 2024.
\url{https://arxiv.org/abs/2410.08129v2}.
\bibitem{threedgut}
Q. Wu, J. Martinez Esturo, A. Mirzaei, N. Moenne-Loccoz, and Z. Gojcic.
\emph{3DGUT: Enabling Distorted Cameras and Secondary Rays in Gaussian
Splatting}. 2024. \url{https://arxiv.org/abs/2412.12507v2}.
\bibitem{oav}
B. Miller, H. Chen, A. Lai, and I. Gkioulekas.
\emph{Objects as volumes: A stochastic geometry view of opaque solids}.
2024. \url{https://arxiv.org/abs/2312.15406v2}.
\bibitem{siegel}
J. W. Siegel. \emph{Optimal Approximation of Zonoids and Uniform
Approximation by Shallow Neural Networks}.
2023. \url{https://arxiv.org/abs/2307.15285v2}.
\bibitem{torres}
G. D. Torres. \emph{On Finding the Closest Zonotope to a Polytope in
Hausdorff Distance}. 2024. \url{https://arxiv.org/abs/2407.13125v1}.
\bibitem{mpmathiv}
mpmath developers. \emph{Contexts}, mpmath 1.3.0 documentation,
section on experimental interval arithmetic.
\url{https://raw.githubusercontent.com/mpmath/mpmath/1.3.0/docs/contexts.rst}.
\bibitem{flicker}
W. Ou, Z. Wu, Y. Zhang, D. Wu, F. Z. Hong, and C. P. Yue.
\emph{FLICKER: A Fine-Grained Contribution-Aware Accelerator for
Real-Time 3D Gaussian Splatting}. 2026, Section IV-C.
\url{https://arxiv.org/abs/2603.01158v1}.
\bibitem{cudamath}
NVIDIA. \emph{CUDA C++ Programming Guide}, version 12.8,
Section 13.1 (Mathematical Functions) and double-precision intrinsic table.
\url{https://docs.nvidia.com/cuda/archive/12.8.1/cuda-c-programming-guide/index.html}.
\bibitem{rvar}
T. Fissler and J. F. Ziegel.
\emph{On the elicitability of range value at risk}.
Statistics \& Risk Modeling, 38(1--2):25--46, 2021.
Preprint: \emph{Evaluating Range Value at Risk Forecasts},
\url{https://arxiv.org/abs/1902.04489v3}.
\bibitem{twodgs}
Binbin Huang et al.
\emph{2D Gaussian Splatting for Geometrically Accurate Radiance Fields}.
ACM Transactions on Graphics, 2024.
\url{https://arxiv.org/abs/2403.17888}.

\bibitem{gggs}
B. Zhang, C. Jiang, H. Li, S. Shen, and P. Tan.
\emph{Geometry-Grounded Gaussian Splatting}. 2026.
\url{https://arxiv.org/abs/2601.17835}.
\bibitem{trianglesplat}
J. Held et al.
\emph{Triangle Splatting for Real-Time Radiance Field Rendering}.
2025. \url{https://arxiv.org/abs/2505.19175v1}.
\bibitem{wrapping}
D. Gomez, A. Guedon, N. Maruani, B. Gong, and M. Ovsjanikov.
\emph{From Blobs to Spokes: High-Fidelity Surface Reconstruction via Oriented
Gaussians}. 2026. \url{https://arxiv.org/abs/2604.07337}.
\bibitem{onestep}
J. Bolte, E. Pauwels, and S. Vaiter.
\emph{One-step differentiation of iterative algorithms}. NeurIPS, 2023.
\url{https://arxiv.org/abs/2305.13768}.
\bibitem{brentdocs}
SciPy developers. \emph{scipy.optimize.brentq}, official documentation.
\url{https://docs.scipy.org/doc/scipy/reference/generated/scipy.optimize.brentq.html}.
\bibitem{dlmf}
NIST Digital Library of Mathematical Functions. Section 3.8, Equation 3.8.12.
\url{https://dlmf.nist.gov/3.8.E12}.
\bibitem{skills}
T. Kassis, V. Agarwal, Y. He, D. Patel, and A. M. Brueckner.
\emph{Scientific Agent Skills: A Library of Procedural Knowledge for Research
Agents}. 2026. \url{https://doi.org/10.48550/arXiv.2609.00065}.
\bibitem{gggscode}
HKUST-SAIL. \emph{Geometry-Grounded-Gaussian-Splatting}, official source,
\texttt{cuda\_rasterizer/render\_forward.cu} and
\texttt{cuda\_rasterizer/render\_backward.cu}, main branch inspected
2026-10-12; not a pinned upstream commit.
\url{https://github.com/HKUST-SAIL/Geometry-Grounded-Gaussian-Splatting}.
\bibitem{fastergs}
F. Hahlbohm, L. Franke, M. Eisemann, and M. Magnor.
\emph{Faster-GS: Analyzing and Improving Gaussian Splatting Optimization}.
2026. Appendix C.2. \url{https://arxiv.org/abs/2602.09999}.
\end{thebibliography}
\end{document}
